\IfFileExists{revtex4-2.cls}
  {\documentclass[aps,prxquantum,preprint,groupedaddress,nofootinbib]{revtex4-2}}
  {\documentclass[11pt]{article}}

\usepackage[T1]{fontenc}
\usepackage{lmodern}
\usepackage{microtype}
\usepackage{amsmath,amssymb,amsfonts,amsthm}
\usepackage{booktabs}
\usepackage{array}
\usepackage{tabularx}
\usepackage{graphicx}
\usepackage{xcolor}
\usepackage{url}
\usepackage[hidelinks]{hyperref}
\hypersetup{
  pdftitle={Predictable Sparse and Spectral Contrast Monitoring of Real QEC Syndromes:
  A Predeclared Benchmark with a Disclosed Post-Detector Repair},
  pdfauthor={Rongzhou (Lachlan) Chen},
  pdfsubject={Predeclared dual-provenance benchmark with a disclosed post-detector repair}
}

\makeatletter
\newif\ifrunrevtex
\@ifclassloaded{revtex4-2}{\runrevtextrue}{\runrevtexfalse}
\makeatother
\ifrunrevtex\else
  \usepackage[numbers,sort&compress]{natbib}
  \usepackage[margin=1in]{geometry}
\fi

\definecolor{pending}{RGB}{166,35,35}
\newcommand{\RunSixPending}[1]{%
  \textcolor{pending}{\nolinkurl{TBD(#1)}}%
}
\newif\ifrunresultbundle
\IfFileExists{generated/publication_bundle_manifest.json}
  {\runresultbundletrue}
  {\runresultbundlefalse}
\ifrunresultbundle
  % Generated by publication/run6/analysis/extract_results.py.
% Override this directory before input if the bundle is not at generated/.
\providecommand{\RunSixGeneratedDir}{generated}
\newcommand{\RunSixVerifiedClaim}{No demonstrated S-PACE algorithmic advantage.
}
\newcommand{\RunSixVerifiedAbstractFinding}{The composite missed the primary Google event window and captured 0 of 31 primary decoder mismatches at the top-20 budget, versus 9 for DFR and 0 for online logistic. Its PNNL macro restricted delay was 0.8046, versus 0.7899 for DFR and 0.8480 for online logistic; the required retention gate failed.}
\newcommand{\RunSixVerifiedOutcomeNarrative}{The Google gate failed: the composite had 0 pre-event alerts, missed the primary event window, and captured 0/31 primary mismatches at top 20, compared with 9/31 for DFR and 0/31 for online logistic. The PNNL gate failed: composite-minus-DFR macro delay was +0.0147, whereas composite-minus-logistic was -0.0434. Retention required strict improvement over both comparators.}
\newcommand{\RunSixVerifiedGateTable}{\noindent\resizebox{\linewidth}{!}{%
\begin{tabular}{p{0.55\linewidth}p{0.27\linewidth}c}
\toprule
Locked condition & Status in protocol & Result \\
\midrule
S-PACE alerts in the primary event window & Google primary gate & Fail \\
S-PACE has at most nine pre-event alerts & Google primary gate & Pass \\
Detector and resource artifacts frozen before labels & Google primary gate & Pass \\
Top-20 capture is at least DFR plus one & Google primary gate & Fail \\
Top-20 capture is at least online logistic plus one & Google primary gate & Fail \\
All windows and locked uncertainty are reported & Google primary gate & Pass \\
All methods receive the same records and labels & Google primary gate & Pass \\
Google aggregate gate & Primary empirical arm & Fail \\
PNNL retention gate & Required auxiliary arm & Fail \\
Overall locked dual-provenance gate & Google AND PNNL & Fail \\
\bottomrule
\end{tabular}
}
\par
}
\newcommand{\RunSixVerifiedGoogleEventTable}{\noindent\resizebox{\linewidth}{!}{%
\begin{tabular}{llrp{0.22\linewidth}rrrrr}
\toprule
Method & Preregistered role & Pre-event & Pre-event shots & Primary & First shot & Role & Narrow & Wide \\
\midrule
DFR & Predesignated comparator & 3 & 44266, 55036, 57712 & Pass & 57750 & 25 & Pass & Pass \\
Within-shot Page CUSUM & Mandatory contextual control & 21 & 41301, 41317, 41727, 42795, 42864, 43105, 44243, 44820, 44921, 45234, 47735, 49107, 52234, 54368, 54491, 54650, 54669, 55210, 55541, 56698, 56904 & Fail & --- & --- & Fail & Fail \\
Diagonal likelihood & Mandatory contextual control & 8 & 41919, 44266, 45673, 52599, 55036, 55298, 57454, 57712 & Pass & 57750 & 25 & Pass & Pass \\
Hotelling & Mandatory contextual control & 24 & 40509, 42264, 42608, 43002, 43025, 43096, 43776, 45284, 46192, 47906, 48384, 48973, 49031, 49318, 51551, 51967, 51973, 53902, 53906, 54378, 54700, 55020, 57651, 57745 & Pass & 57750 & 25 & Pass & Pass \\
Online logistic & Predesignated comparator & 13 & 40177, 40190, 40220, 40243, 40335, 40408, 40918, 41303, 41579, 47791, 54012, 54619, 57338 & Fail & --- & --- & Fail & Fail \\
S-PACE sparse & Preregistered component & 0 & none & Fail & --- & --- & Fail & Fail \\
S-PACE spectral & Preregistered component & 6 & 44907, 50797, 52804, 54012, 54748, 57646 & Pass & 57778 & 9 & Pass & Pass \\
S-PACE composite & Primary proposed method & 0 & none & Fail & --- & --- & Fail & Fail \\
\bottomrule
\end{tabular}
}
\par
}
\newcommand{\RunSixVerifiedGoogleThresholdFrontierTable}{\noindent\resizebox{\linewidth}{!}{%
\begin{tabular}{llrrrrl}
\toprule
Method & Preregistered role & Primary threshold & Validation alerts/budget & Zero-alert threshold & Zero-alert count & Complete frontier \\
\midrule
DFR & Predesignated comparator & 0.541667 & 0/2 & 0.541667 & 0 & Hash-verified \\
Within-shot Page CUSUM & Mandatory contextual control & 35.500000 & 2/2 & 39.540000 & 0 & Hash-verified \\
Diagonal likelihood & Mandatory contextual control & 0.114604 & 2/2 & 0.120120 & 0 & Hash-verified \\
Hotelling & Mandatory contextual control & 2431.422796 & 2/2 & 2505.799438 & 0 & Hash-verified \\
Online logistic & Predesignated comparator & 0.376278 & 2/2 & 0.409261 & 0 & Hash-verified \\
S-PACE sparse & Preregistered component & 1.000000 & 0/2 & 1.000000 & 0 & Hash-verified \\
S-PACE spectral & Preregistered component & 0.812654 & 2/2 & 0.862347 & 0 & Hash-verified \\
S-PACE composite & Primary proposed method & 1.000000 & 0/2 & 1.000000 & 0 & Hash-verified \\
\bottomrule
\end{tabular}
}
\par
}
\newcommand{\RunSixVerifiedGoogleRiskBudgetTable}{\noindent\resizebox{\linewidth}{!}{%
\begin{tabular}{lrlp{0.20\linewidth}rrrrr}
\toprule
Label & Budget & Method & Preregistered role & Captured/total [95\%; n] & Recall [95\%; n] & Precision [95\%; n] & Retained rate [95\%; n] & Coverage \\
\midrule
Primary: correlated matching & 2 & DFR & Predesignated comparator & 1/31 1.00 [0.00, 2.00]; n=2000 & 0.03226 [0.00000, 0.16667]; n=1632 & 0.50000 [0.00000, 1.00000]; n=2000 & 0.00150 [0.00000, 0.00470]; n=2000 & 0.99990 \\
Primary: correlated matching & 2 & Within-shot Page CUSUM & Mandatory contextual control & 0/31 0.00 [0.00, 0.00]; n=2000 & 0.00000 [0.00000, 0.00000]; n=1632 & 0.00000 [0.00000, 0.00000]; n=2000 & 0.00155 [0.00000, 0.00480]; n=2000 & 0.99990 \\
Primary: correlated matching & 2 & Diagonal likelihood & Mandatory contextual control & 1/31 1.00 [0.00, 1.00]; n=2000 & 0.03226 [0.00000, 0.14286]; n=1632 & 0.50000 [0.00000, 0.50000]; n=2000 & 0.00150 [0.00000, 0.00480]; n=2000 & 0.99990 \\
Primary: correlated matching & 2 & Hotelling & Mandatory contextual control & 0/31 0.00 [0.00, 1.00]; n=2000 & 0.00000 [0.00000, 0.07692]; n=1632 & 0.00000 [0.00000, 0.50000]; n=2000 & 0.00155 [0.00000, 0.00480]; n=2000 & 0.99990 \\
Primary: correlated matching & 2 & Online logistic & Predesignated comparator & 0/31 0.00 [0.00, 0.00]; n=2000 & 0.00000 [0.00000, 0.00000]; n=1632 & 0.00000 [0.00000, 0.00000]; n=2000 & 0.00155 [0.00000, 0.00480]; n=2000 & 0.99990 \\
Primary: correlated matching & 2 & S-PACE sparse & Preregistered component & 0/31 0.00 [0.00, 0.00]; n=2000 & 0.00000 [0.00000, 0.00000]; n=1632 & 0.00000 [0.00000, 0.00000]; n=2000 & 0.00155 [0.00000, 0.00480]; n=2000 & 0.99990 \\
Primary: correlated matching & 2 & S-PACE spectral & Preregistered component & 0/31 0.00 [0.00, 1.00]; n=2000 & 0.00000 [0.00000, 0.04037]; n=1632 & 0.00000 [0.00000, 0.50000]; n=2000 & 0.00155 [0.00000, 0.00480]; n=2000 & 0.99990 \\
Primary: correlated matching & 2 & S-PACE composite & Primary proposed method & 0/31 0.00 [0.00, 0.00]; n=2000 & 0.00000 [0.00000, 0.00000]; n=1632 & 0.00000 [0.00000, 0.00000]; n=2000 & 0.00155 [0.00000, 0.00480]; n=2000 & 0.99990 \\
Primary: correlated matching & 20 & DFR & Predesignated comparator & 9/31 9.00 [0.00, 12.00]; n=2000 & 0.29032 [0.00000, 0.50000]; n=1632 & 0.45000 [0.00000, 0.60000]; n=2000 & 0.00110 [0.00000, 0.00445]; n=2000 & 0.99900 \\
Primary: correlated matching & 20 & Within-shot Page CUSUM & Mandatory contextual control & 0/31 0.00 [0.00, 0.00]; n=2000 & 0.00000 [0.00000, 0.00000]; n=1632 & 0.00000 [0.00000, 0.00000]; n=2000 & 0.00155 [0.00000, 0.00480]; n=2000 & 0.99900 \\
Primary: correlated matching & 20 & Diagonal likelihood & Mandatory contextual control & 9/31 9.00 [0.00, 10.00]; n=2000 & 0.29032 [0.00000, 0.50000]; n=1632 & 0.45000 [0.00000, 0.50000]; n=2000 & 0.00110 [0.00000, 0.00460]; n=2000 & 0.99900 \\
Primary: correlated matching & 20 & Hotelling & Mandatory contextual control & 6/31 6.00 [0.00, 8.00]; n=2000 & 0.19355 [0.00000, 0.46154]; n=1632 & 0.30000 [0.00000, 0.40000]; n=2000 & 0.00125 [0.00000, 0.00465]; n=2000 & 0.99900 \\
Primary: correlated matching & 20 & Online logistic & Predesignated comparator & 0/31 0.00 [0.00, 0.00]; n=2000 & 0.00000 [0.00000, 0.00000]; n=1632 & 0.00000 [0.00000, 0.00000]; n=2000 & 0.00155 [0.00000, 0.00480]; n=2000 & 0.99900 \\
Primary: correlated matching & 20 & S-PACE sparse & Preregistered component & 0/31 0.00 [0.00, 0.00]; n=2000 & 0.00000 [0.00000, 0.00000]; n=1632 & 0.00000 [0.00000, 0.00000]; n=2000 & 0.00155 [0.00000, 0.00480]; n=2000 & 0.99900 \\
Primary: correlated matching & 20 & S-PACE spectral & Preregistered component & 3/31 3.00 [0.00, 8.00]; n=2000 & 0.09677 [0.00000, 0.25000]; n=1632 & 0.15000 [0.00000, 0.40000]; n=2000 & 0.00140 [0.00000, 0.00456]; n=2000 & 0.99900 \\
Primary: correlated matching & 20 & S-PACE composite & Primary proposed method & 0/31 0.00 [0.00, 0.00]; n=2000 & 0.00000 [0.00000, 0.00000]; n=1632 & 0.00000 [0.00000, 0.00000]; n=2000 & 0.00155 [0.00000, 0.00480]; n=2000 & 0.99900 \\
Primary: correlated matching & 200 & DFR & Predesignated comparator & 25/31 25.00 [0.00, 72.00]; n=2000 & 0.80645 [0.00000, 0.90107]; n=1632 & 0.12500 [0.00000, 0.36000]; n=2000 & 0.00030 [0.00000, 0.00126]; n=2000 & 0.99000 \\
Primary: correlated matching & 200 & Within-shot Page CUSUM & Mandatory contextual control & 3/31 3.00 [0.00, 9.00]; n=2000 & 0.09677 [0.00000, 0.28571]; n=1632 & 0.01500 [0.00000, 0.04500]; n=2000 & 0.00141 [0.00000, 0.00439]; n=2000 & 0.99000 \\
Primary: correlated matching & 200 & Diagonal likelihood & Mandatory contextual control & 28/31 28.00 [0.00, 68.00]; n=2000 & 0.90323 [0.00000, 0.95833]; n=1632 & 0.14000 [0.00000, 0.34000]; n=2000 & 0.00015 [0.00000, 0.00162]; n=2000 & 0.99000 \\
Primary: correlated matching & 200 & Hotelling & Mandatory contextual control & 27/31 27.00 [0.00, 64.00]; n=2000 & 0.87097 [0.00000, 1.00000]; n=1632 & 0.13500 [0.00000, 0.32000]; n=2000 & 0.00020 [0.00000, 0.00157]; n=2000 & 0.99000 \\
Primary: correlated matching & 200 & Online logistic & Predesignated comparator & 1/31 1.00 [0.00, 4.00]; n=2000 & 0.03226 [0.00000, 1.00000]; n=1632 & 0.00500 [0.00000, 0.02000]; n=2000 & 0.00152 [0.00000, 0.00480]; n=2000 & 0.99000 \\
Primary: correlated matching & 200 & S-PACE sparse & Preregistered component & 0/31 0.00 [0.00, 0.00]; n=2000 & 0.00000 [0.00000, 0.00000]; n=1632 & 0.00000 [0.00000, 0.00000]; n=2000 & 0.00157 [0.00000, 0.00485]; n=2000 & 0.99000 \\
Primary: correlated matching & 200 & S-PACE spectral & Preregistered component & 17/31 17.00 [0.00, 44.00]; n=2000 & 0.54839 [0.00000, 0.81939]; n=1632 & 0.08500 [0.00000, 0.22000]; n=2000 & 0.00071 [0.00000, 0.00268]; n=2000 & 0.99000 \\
Primary: correlated matching & 200 & S-PACE composite & Primary proposed method & 0/31 0.00 [0.00, 0.00]; n=2000 & 0.00000 [0.00000, 0.00000]; n=1632 & 0.00000 [0.00000, 0.00000]; n=2000 & 0.00157 [0.00000, 0.00485]; n=2000 & 0.99000 \\
Secondary: PyMatching & 2 & DFR & Predesignated comparator & 1/30 1.00 [0.00, 2.00]; n=2000 & 0.03333 [0.00000, 0.16667]; n=1632 & 0.50000 [0.00000, 1.00000]; n=2000 & 0.00145 [0.00000, 0.00450]; n=2000 & 0.99990 \\
Secondary: PyMatching & 2 & Within-shot Page CUSUM & Mandatory contextual control & 0/30 0.00 [0.00, 0.00]; n=2000 & 0.00000 [0.00000, 0.00000]; n=1632 & 0.00000 [0.00000, 0.00000]; n=2000 & 0.00150 [0.00000, 0.00460]; n=2000 & 0.99990 \\
Secondary: PyMatching & 2 & Diagonal likelihood & Mandatory contextual control & 1/30 1.00 [0.00, 1.00]; n=2000 & 0.03333 [0.00000, 0.14286]; n=1632 & 0.50000 [0.00000, 0.50000]; n=2000 & 0.00145 [0.00000, 0.00460]; n=2000 & 0.99990 \\
Secondary: PyMatching & 2 & Hotelling & Mandatory contextual control & 0/30 0.00 [0.00, 1.00]; n=2000 & 0.00000 [0.00000, 0.07692]; n=1632 & 0.00000 [0.00000, 0.50000]; n=2000 & 0.00150 [0.00000, 0.00460]; n=2000 & 0.99990 \\
Secondary: PyMatching & 2 & Online logistic & Predesignated comparator & 0/30 0.00 [0.00, 0.00]; n=2000 & 0.00000 [0.00000, 0.00000]; n=1632 & 0.00000 [0.00000, 0.00000]; n=2000 & 0.00150 [0.00000, 0.00460]; n=2000 & 0.99990 \\
Secondary: PyMatching & 2 & S-PACE sparse & Preregistered component & 0/30 0.00 [0.00, 0.00]; n=2000 & 0.00000 [0.00000, 0.00000]; n=1632 & 0.00000 [0.00000, 0.00000]; n=2000 & 0.00150 [0.00000, 0.00460]; n=2000 & 0.99990 \\
Secondary: PyMatching & 2 & S-PACE spectral & Preregistered component & 0/30 0.00 [0.00, 1.00]; n=2000 & 0.00000 [0.00000, 0.04207]; n=1632 & 0.00000 [0.00000, 0.50000]; n=2000 & 0.00150 [0.00000, 0.00460]; n=2000 & 0.99990 \\
Secondary: PyMatching & 2 & S-PACE composite & Primary proposed method & 0/30 0.00 [0.00, 0.00]; n=2000 & 0.00000 [0.00000, 0.00000]; n=1632 & 0.00000 [0.00000, 0.00000]; n=2000 & 0.00150 [0.00000, 0.00460]; n=2000 & 0.99990 \\
Secondary: PyMatching & 20 & DFR & Predesignated comparator & 9/30 9.00 [0.00, 12.00]; n=2000 & 0.30000 [0.00000, 0.50000]; n=1632 & 0.45000 [0.00000, 0.60000]; n=2000 & 0.00105 [0.00000, 0.00426]; n=2000 & 0.99900 \\
Secondary: PyMatching & 20 & Within-shot Page CUSUM & Mandatory contextual control & 0/30 0.00 [0.00, 0.00]; n=2000 & 0.00000 [0.00000, 0.00000]; n=1632 & 0.00000 [0.00000, 0.00000]; n=2000 & 0.00150 [0.00000, 0.00461]; n=2000 & 0.99900 \\
Secondary: PyMatching & 20 & Diagonal likelihood & Mandatory contextual control & 9/30 9.00 [0.00, 9.00]; n=2000 & 0.30000 [0.00000, 0.50000]; n=1632 & 0.45000 [0.00000, 0.45000]; n=2000 & 0.00105 [0.00000, 0.00441]; n=2000 & 0.99900 \\
Secondary: PyMatching & 20 & Hotelling & Mandatory contextual control & 6/30 6.00 [0.00, 8.00]; n=2000 & 0.20000 [0.00000, 0.46154]; n=1632 & 0.30000 [0.00000, 0.40000]; n=2000 & 0.00120 [0.00000, 0.00446]; n=2000 & 0.99900 \\
Secondary: PyMatching & 20 & Online logistic & Predesignated comparator & 0/30 0.00 [0.00, 0.00]; n=2000 & 0.00000 [0.00000, 0.00000]; n=1632 & 0.00000 [0.00000, 0.00000]; n=2000 & 0.00150 [0.00000, 0.00461]; n=2000 & 0.99900 \\
Secondary: PyMatching & 20 & S-PACE sparse & Preregistered component & 0/30 0.00 [0.00, 0.00]; n=2000 & 0.00000 [0.00000, 0.00000]; n=1632 & 0.00000 [0.00000, 0.00000]; n=2000 & 0.00150 [0.00000, 0.00461]; n=2000 & 0.99900 \\
Secondary: PyMatching & 20 & S-PACE spectral & Preregistered component & 3/30 3.00 [0.00, 7.00]; n=2000 & 0.10000 [0.00000, 0.25000]; n=1632 & 0.15000 [0.00000, 0.35000]; n=2000 & 0.00135 [0.00000, 0.00440]; n=2000 & 0.99900 \\
Secondary: PyMatching & 20 & S-PACE composite & Primary proposed method & 0/30 0.00 [0.00, 0.00]; n=2000 & 0.00000 [0.00000, 0.00000]; n=1632 & 0.00000 [0.00000, 0.00000]; n=2000 & 0.00150 [0.00000, 0.00461]; n=2000 & 0.99900 \\
Secondary: PyMatching & 200 & DFR & Predesignated comparator & 24/30 24.00 [0.00, 68.00]; n=2000 & 0.80000 [0.00000, 0.89733]; n=1632 & 0.12000 [0.00000, 0.34000]; n=2000 & 0.00030 [0.00000, 0.00126]; n=2000 & 0.99000 \\
Secondary: PyMatching & 200 & Within-shot Page CUSUM & Mandatory contextual control & 3/30 3.00 [0.00, 9.00]; n=2000 & 0.10000 [0.00000, 0.28571]; n=1632 & 0.01500 [0.00000, 0.04500]; n=2000 & 0.00136 [0.00000, 0.00419]; n=2000 & 0.99000 \\
Secondary: PyMatching & 200 & Diagonal likelihood & Mandatory contextual control & 27/30 27.00 [0.00, 64.00]; n=2000 & 0.90000 [0.00000, 0.95652]; n=1632 & 0.13500 [0.00000, 0.32000]; n=2000 & 0.00015 [0.00000, 0.00162]; n=2000 & 0.99000 \\
Secondary: PyMatching & 200 & Hotelling & Mandatory contextual control & 26/30 26.00 [0.00, 60.00]; n=2000 & 0.86667 [0.00000, 1.00000]; n=1632 & 0.13000 [0.00000, 0.30000]; n=2000 & 0.00020 [0.00000, 0.00157]; n=2000 & 0.99000 \\
Secondary: PyMatching & 200 & Online logistic & Predesignated comparator & 1/30 1.00 [0.00, 4.00]; n=2000 & 0.03333 [0.00000, 1.00000]; n=1632 & 0.00500 [0.00000, 0.02000]; n=2000 & 0.00146 [0.00000, 0.00465]; n=2000 & 0.99000 \\
Secondary: PyMatching & 200 & S-PACE sparse & Preregistered component & 0/30 0.00 [0.00, 0.00]; n=2000 & 0.00000 [0.00000, 0.00000]; n=1632 & 0.00000 [0.00000, 0.00000]; n=2000 & 0.00152 [0.00000, 0.00465]; n=2000 & 0.99000 \\
Secondary: PyMatching & 200 & S-PACE spectral & Preregistered component & 16/30 16.00 [0.00, 41.02]; n=2000 & 0.53333 [0.00000, 0.81250]; n=1632 & 0.08000 [0.00000, 0.20512]; n=2000 & 0.00071 [0.00000, 0.00268]; n=2000 & 0.99000 \\
Secondary: PyMatching & 200 & S-PACE composite & Primary proposed method & 0/30 0.00 [0.00, 0.00]; n=2000 & 0.00000 [0.00000, 0.00000]; n=1632 & 0.00000 [0.00000, 0.00000]; n=2000 & 0.00152 [0.00000, 0.00465]; n=2000 & 0.99000 \\
\bottomrule
\end{tabular}
}
\par
}
\newcommand{\RunSixVerifiedGoogleUncertaintyTable}{\noindent\resizebox{\linewidth}{!}{%
\begin{tabular}{lrlrrrrl}
\toprule
Label & Budget & Locked contrast & Captured difference [95\%; n] & Recall difference [95\%; n] & Precision difference [95\%; n] & Retained-rate difference [95\%; n] & Gate relevance \\
\midrule
Primary: correlated matching & 2 & S-PACE minus DFR & -1.00 [-2.00, 0.00]; n=2000 & -0.03226 [-0.16667, 0.00000]; n=1632 & -0.50000 [-1.00000, 0.00000]; n=2000 & 0.00005 [0.00000, 0.00010]; n=2000 & Descriptive \\
Primary: correlated matching & 20 & S-PACE minus DFR & -9.00 [-12.00, 0.00]; n=2000 & -0.29032 [-0.50000, 0.00000]; n=1632 & -0.45000 [-0.60000, 0.00000]; n=2000 & 0.00045 [0.00000, 0.00060]; n=2000 & Gate capture contrast \\
Primary: correlated matching & 200 & S-PACE minus DFR & -25.00 [-72.00, 0.00]; n=2000 & -0.80645 [-0.90107, 0.00000]; n=1632 & -0.12500 [-0.36000, 0.00000]; n=2000 & 0.00126 [0.00000, 0.00364]; n=2000 & Descriptive \\
Primary: correlated matching & 2 & S-PACE minus Online logistic & 0.00 [0.00, 0.00]; n=2000 & 0.00000 [0.00000, 0.00000]; n=1632 & 0.00000 [0.00000, 0.00000]; n=2000 & 0.00000 [0.00000, 0.00000]; n=2000 & Descriptive \\
Primary: correlated matching & 20 & S-PACE minus Online logistic & 0.00 [0.00, 0.00]; n=2000 & 0.00000 [0.00000, 0.00000]; n=1632 & 0.00000 [0.00000, 0.00000]; n=2000 & 0.00000 [0.00000, 0.00000]; n=2000 & Gate capture contrast \\
Primary: correlated matching & 200 & S-PACE minus Online logistic & -1.00 [-4.00, 0.00]; n=2000 & -0.03226 [-1.00000, 0.00000]; n=1632 & -0.00500 [-0.02000, 0.00000]; n=2000 & 0.00005 [0.00000, 0.00020]; n=2000 & Descriptive \\
Secondary: PyMatching & 2 & S-PACE minus DFR & -1.00 [-2.00, 0.00]; n=2000 & -0.03333 [-0.16667, 0.00000]; n=1632 & -0.50000 [-1.00000, 0.00000]; n=2000 & 0.00005 [0.00000, 0.00010]; n=2000 & Descriptive \\
Secondary: PyMatching & 20 & S-PACE minus DFR & -9.00 [-12.00, 0.00]; n=2000 & -0.30000 [-0.50000, 0.00000]; n=1632 & -0.45000 [-0.60000, 0.00000]; n=2000 & 0.00045 [0.00000, 0.00060]; n=2000 & Descriptive \\
Secondary: PyMatching & 200 & S-PACE minus DFR & -24.00 [-68.00, 0.00]; n=2000 & -0.80000 [-0.89733, 0.00000]; n=1632 & -0.12000 [-0.34000, 0.00000]; n=2000 & 0.00121 [0.00000, 0.00343]; n=2000 & Descriptive \\
Secondary: PyMatching & 2 & S-PACE minus Online logistic & 0.00 [0.00, 0.00]; n=2000 & 0.00000 [0.00000, 0.00000]; n=1632 & 0.00000 [0.00000, 0.00000]; n=2000 & 0.00000 [0.00000, 0.00000]; n=2000 & Descriptive \\
Secondary: PyMatching & 20 & S-PACE minus Online logistic & 0.00 [0.00, 0.00]; n=2000 & 0.00000 [0.00000, 0.00000]; n=1632 & 0.00000 [0.00000, 0.00000]; n=2000 & 0.00000 [0.00000, 0.00000]; n=2000 & Descriptive \\
Secondary: PyMatching & 200 & S-PACE minus Online logistic & -1.00 [-4.00, 0.00]; n=2000 & -0.03333 [-1.00000, 0.00000]; n=1632 & -0.00500 [-0.02000, 0.00000]; n=2000 & 0.00005 [0.00000, 0.00020]; n=2000 & Descriptive \\
\bottomrule
\end{tabular}
}
\par
}
\newcommand{\RunSixVerifiedGoogleThresholdBootstrapTable}{\noindent\resizebox{\linewidth}{!}{%
\begin{tabular}{lrrrrrr}
\toprule
Method & Frozen threshold & Selected threshold (2.5/50/97.5\%) & Selected-alert histogram & Alerts at frozen threshold (2.5/50/97.5\%) & Frozen-alert histogram & Zero-alert threshold (2.5/50/97.5\%) \\
\midrule
DFR & 0.541667 & 0.500000 / 0.541667 / 0.541667 & 0:1301, 1:265, 2:434 & 0.00 / 0.00 / 0.00 & 0:2000 & 0.500000 / 0.541667 / 0.541667 \\
Within-shot Page CUSUM & 35.500000 & 31.630000 / 35.500000 / 39.540000 & 0:275, 1:479, 2:1246 & 0.00 / 2.00 / 5.00 & 0:226, 1:533, 2:582, 3:375, 4:194, 5:63, 6:21, 7:3, 8:3 & 33.570000 / 39.540000 / 39.540000 \\
Diagonal likelihood & 0.114604 & 0.110576 / 0.114604 / 0.120120 & 0:272, 1:490, 2:1238 & 0.00 / 2.00 / 5.00 & 0:222, 1:531, 2:557, 3:366, 4:225, 5:78, 6:16, 7:4, 9:1 & 0.114293 / 0.120120 / 0.120120 \\
Hotelling & 2431.422796 & 2380.478942 / 2431.422796 / 2505.799438 & 0:254, 1:490, 2:1256 & 0.00 / 2.00 / 5.00 & 0:241, 1:569, 2:537, 3:377, 4:197, 5:53, 6:22, 7:4 & 2425.313379 / 2505.799438 / 2505.799438 \\
Online logistic & 0.376278 & 0.352765 / 0.376278 / 0.409261 & 0:247, 1:393, 2:1360 & 0.00 / 2.00 / 5.00 & 0:268, 1:544, 2:542, 3:360, 4:193, 5:66, 6:20, 7:4, 8:3 & 0.362359 / 0.409261 / 0.409261 \\
S-PACE sparse & 1.000000 & 1.000000 / 1.000000 / 1.000000 & 0:2000 & 0.00 / 0.00 / 0.00 & 0:2000 & 1.000000 / 1.000000 / 1.000000 \\
S-PACE spectral & 0.812654 & 0.802795 / 0.812654 / 0.862347 & 0:238, 1:519, 2:1243 & 0.00 / 2.00 / 5.00 & 0:257, 1:552, 2:541, 3:387, 4:176, 5:65, 6:15, 7:6, 8:1 & 0.805416 / 0.862347 / 0.862347 \\
S-PACE composite & 1.000000 & 1.000000 / 1.000000 / 1.000000 & 0:2000 & 0.00 / 0.00 / 0.00 & 0:2000 & 1.000000 / 1.000000 / 1.000000 \\
\bottomrule
\end{tabular}
}
\par
}
\newcommand{\RunSixVerifiedGoogleRandomizationTable}{\noindent\resizebox{\linewidth}{!}{%
\begin{tabular}{llrrrr}
\toprule
Method/statistic & Proper-prior quantity & Count & Replicates & Fraction & 95\% interval \\
\midrule
DFR & Proper-prior e-process >= 100 & 0 & 256 & 0.0000 & Not preregistered \\
Diagonal likelihood & Proper-prior e-process >= 100 & 0 & 256 & 0.0000 & Not preregistered \\
Online logistic & Proper-prior e-process >= 100 & 0 & 256 & 0.0000 & Not preregistered \\
S-PACE sparse & Proper-prior e-process >= 100 & 1 & 256 & 0.0039 & Not preregistered \\
S-PACE spectral & Proper-prior e-process >= 100 & 0 & 256 & 0.0000 & Not preregistered \\
S-PACE composite & Proper-prior e-process >= 100 & 1 & 256 & 0.0039 & [0.0001, 0.0216] \\
Any-family crossing at 600 & Proper-prior family diagnostic & 0 & 256 & --- & --- \\
\bottomrule
\end{tabular}
}
\par
}
\newcommand{\RunSixVerifiedPNNLMacroTable}{\noindent\resizebox{\linewidth}{!}{%
\begin{tabular}{llrrr}
\toprule
Method & Preregistered role & Pre-FA states & Miss states & Macro delay \\
\midrule
DFR & Predesignated comparator & 0 & 16 & 0.7899 \\
Online logistic & Predesignated comparator & 0 & 16 & 0.8480 \\
S-PACE sparse & Preregistered component & 0 & 16 & 0.8019 \\
S-PACE spectral & Preregistered component & 0 & 15 & 0.8144 \\
S-PACE composite & Primary auxiliary target & 0 & 16 & 0.8046 \\
\bottomrule
\end{tabular}
}
\par
}
\newcommand{\RunSixVerifiedPNNLStateResultsTable}{\noindent\resizebox{\linewidth}{!}{%
\begin{tabular}{llrrrrrr}
\toprule
Plot ID & Cohort ID & State & DFR & Online logistic & S-PACE sparse & S-PACE spectral & S-PACE composite \\
\midrule
C01 & pgh\_d3\_r1\_x\_18\_11\_12\_13\_14 & 0 & 0/1/1.0000 & 0/1/1.0000 & 0/1/1.0000 & 0/1/1.0000 & 0/1/1.0000 \\
C01 & pgh\_d3\_r1\_x\_18\_11\_12\_13\_14 & 1 & 0/1/1.0000 & 0/1/1.0000 & 0/1/1.0000 & 0/1/1.0000 & 0/1/1.0000 \\
C02 & pgh\_d3\_r1\_x\_111\_98\_91\_92\_93 & 0 & 0/1/1.0000 & 0/1/1.0000 & 0/1/1.0000 & 0/1/1.0000 & 0/1/1.0000 \\
C02 & pgh\_d3\_r1\_x\_111\_98\_91\_92\_93 & 1 & 0/1/1.0000 & 0/1/1.0000 & 0/1/1.0000 & 0/1/1.0000 & 0/1/1.0000 \\
C03 & pgh\_d3\_r3\_x\_14\_13\_12\_11\_18 & 0 & 0/1/1.0000 & 0/1/1.0000 & 0/1/1.0000 & 0/1/1.0000 & 0/1/1.0000 \\
C03 & pgh\_d3\_r3\_x\_14\_13\_12\_11\_18 & 1 & 0/1/1.0000 & 0/1/1.0000 & 0/1/1.0000 & 0/1/1.0000 & 0/1/1.0000 \\
C04 & pgh\_d3\_r3\_x\_115\_114\_113\_119\_133 & 0 & 0/1/1.0000 & 0/1/1.0000 & 0/1/1.0000 & 0/1/1.0000 & 0/1/1.0000 \\
C04 & pgh\_d3\_r3\_x\_115\_114\_113\_119\_133 & 1 & 0/1/1.0000 & 0/1/1.0000 & 0/1/1.0000 & 0/0/0.9487 & 0/1/1.0000 \\
C05 & pgh\_d3\_r3\_z\_151\_152\_153\_154\_155 & 0 & 0/1/1.0000 & 0/1/1.0000 & 0/1/1.0000 & 0/1/1.0000 & 0/1/1.0000 \\
C05 & pgh\_d3\_r3\_z\_151\_152\_153\_154\_155 & 1 & 0/1/1.0000 & 0/1/1.0000 & 0/1/1.0000 & 0/1/1.0000 & 0/1/1.0000 \\
C06 & pgh\_d5\_r1\_x\_18\_11\_12\_13\_14\_15\_19\_35\_34 & 0 & 0/0/0.1921 & 0/0/0.4443 & 0/0/0.1848 & 0/0/0.1716 & 0/0/0.1818 \\
C06 & pgh\_d5\_r1\_x\_18\_11\_12\_13\_14\_15\_19\_35\_34 & 1 & 0/0/0.2155 & 0/0/0.3123 & 0/0/0.3021 & 0/0/0.2947 & 0/0/0.2977 \\
C07 & pgh\_d5\_r1\_z\_18\_11\_12\_13\_14\_15\_19\_35\_34 & 0 & 0/1/1.0000 & 0/1/1.0000 & 0/1/1.0000 & 0/1/1.0000 & 0/1/1.0000 \\
C07 & pgh\_d5\_r1\_z\_18\_11\_12\_13\_14\_15\_19\_35\_34 & 1 & 0/1/1.0000 & 0/1/1.0000 & 0/1/1.0000 & 0/1/1.0000 & 0/1/1.0000 \\
C08 & pgh\_d5\_r5\_x\_34\_35\_19\_15\_14\_13\_12\_11\_18 & 0 & 0/0/0.2097 & 0/0/0.1525 & 0/0/0.1100 & 0/0/0.0953 & 0/0/0.1100 \\
C08 & pgh\_d5\_r5\_x\_34\_35\_19\_15\_14\_13\_12\_11\_18 & 1 & 0/0/0.2390 & 0/0/0.1789 & 0/0/0.1114 & 0/0/0.1496 & 0/0/0.1320 \\
C09 & pgh\_d7\_r3\_z\_80\_81\_82\_83\_84\_85\_77\_65\_66\_67\_57\_47\_46 & 0 & 0/1/1.0000 & 0/1/1.0000 & 0/1/1.0000 & 0/1/1.0000 & 0/1/1.0000 \\
C09 & pgh\_d7\_r3\_z\_80\_81\_82\_83\_84\_85\_77\_65\_66\_67\_57\_47\_46 & 1 & 0/1/1.0000 & 0/1/1.0000 & 0/1/1.0000 & 0/1/1.0000 & 0/1/1.0000 \\
C10 & pgh\_d7\_r7\_z\_138\_151\_150\_149\_148\_147\_146\_145\_144\_143\_142\_141\_140 & 0 & 0/1/1.0000 & 0/1/1.0000 & 0/1/1.0000 & 0/1/1.0000 & 0/1/1.0000 \\
C10 & pgh\_d7\_r7\_z\_138\_151\_150\_149\_148\_147\_146\_145\_144\_143\_142\_141\_140 & 1 & 0/1/1.0000 & 0/1/1.0000 & 0/1/1.0000 & 0/1/1.0000 & 0/1/1.0000 \\
C11 & pgh\_d9\_r5\_x\_153\_154\_155\_139\_135\_134\_133\_119\_113\_114\_115\_99\_95\_94\_93\_92\_91 & 0 & 0/0/0.2522 & 0/0/0.8372 & 0/0/0.2522 & 0/0/0.6804 & 0/0/0.3446 \\
C11 & pgh\_d9\_r5\_x\_153\_154\_155\_139\_135\_134\_133\_119\_113\_114\_115\_99\_95\_94\_93\_92\_91 & 1 & 0/0/0.2698 & 0/0/0.7302 & 0/0/0.6804 & 0/0/0.5762 & 0/0/0.6349 \\
\bottomrule
\end{tabular}
}
\par
}
\newcommand{\RunSixVerifiedPNNLCohortResultsTable}{\noindent\resizebox{\linewidth}{!}{%
\begin{tabular}{llrrrrr}
\toprule
Plot ID & Cohort ID & DFR & Online logistic & S-PACE sparse & S-PACE spectral & S-PACE composite \\
\midrule
C01 & pgh\_d3\_r1\_x\_18\_11\_12\_13\_14 & 0.000/1.000/1.0000 & 0.000/1.000/1.0000 & 0.000/1.000/1.0000 & 0.000/1.000/1.0000 & 0.000/1.000/1.0000 \\
C02 & pgh\_d3\_r1\_x\_111\_98\_91\_92\_93 & 0.000/1.000/1.0000 & 0.000/1.000/1.0000 & 0.000/1.000/1.0000 & 0.000/1.000/1.0000 & 0.000/1.000/1.0000 \\
C03 & pgh\_d3\_r3\_x\_14\_13\_12\_11\_18 & 0.000/1.000/1.0000 & 0.000/1.000/1.0000 & 0.000/1.000/1.0000 & 0.000/1.000/1.0000 & 0.000/1.000/1.0000 \\
C04 & pgh\_d3\_r3\_x\_115\_114\_113\_119\_133 & 0.000/1.000/1.0000 & 0.000/1.000/1.0000 & 0.000/1.000/1.0000 & 0.000/0.500/0.9743 & 0.000/1.000/1.0000 \\
C05 & pgh\_d3\_r3\_z\_151\_152\_153\_154\_155 & 0.000/1.000/1.0000 & 0.000/1.000/1.0000 & 0.000/1.000/1.0000 & 0.000/1.000/1.0000 & 0.000/1.000/1.0000 \\
C06 & pgh\_d5\_r1\_x\_18\_11\_12\_13\_14\_15\_19\_35\_34 & 0.000/0.000/0.2038 & 0.000/0.000/0.3783 & 0.000/0.000/0.2434 & 0.000/0.000/0.2331 & 0.000/0.000/0.2397 \\
C07 & pgh\_d5\_r1\_z\_18\_11\_12\_13\_14\_15\_19\_35\_34 & 0.000/1.000/1.0000 & 0.000/1.000/1.0000 & 0.000/1.000/1.0000 & 0.000/1.000/1.0000 & 0.000/1.000/1.0000 \\
C08 & pgh\_d5\_r5\_x\_34\_35\_19\_15\_14\_13\_12\_11\_18 & 0.000/0.000/0.2243 & 0.000/0.000/0.1657 & 0.000/0.000/0.1107 & 0.000/0.000/0.1224 & 0.000/0.000/0.1210 \\
C09 & pgh\_d7\_r3\_z\_80\_81\_82\_83\_84\_85\_77\_65\_66\_67\_57\_47\_46 & 0.000/1.000/1.0000 & 0.000/1.000/1.0000 & 0.000/1.000/1.0000 & 0.000/1.000/1.0000 & 0.000/1.000/1.0000 \\
C10 & pgh\_d7\_r7\_z\_138\_151\_150\_149\_148\_147\_146\_145\_144\_143\_142\_141\_140 & 0.000/1.000/1.0000 & 0.000/1.000/1.0000 & 0.000/1.000/1.0000 & 0.000/1.000/1.0000 & 0.000/1.000/1.0000 \\
C11 & pgh\_d9\_r5\_x\_153\_154\_155\_139\_135\_134\_133\_119\_113\_114\_115\_99\_95\_94\_93\_92\_91 & 0.000/0.000/0.2610 & 0.000/0.000/0.7837 & 0.000/0.000/0.4663 & 0.000/0.000/0.6283 & 0.000/0.000/0.4897 \\
\bottomrule
\end{tabular}
}
\par
}
\newcommand{\RunSixVerifiedPNNLComparisonTable}{\noindent\resizebox{\linewidth}{!}{%
\begin{tabular}{lrrrrrrr}
\toprule
Comparator & S-PACE minus comparator & Path-bootstrap 95\% interval & Calibration-pair sensitivity 95\% interval & Exact sign-flip p & No-worse pre-FA & Lower delay & Retention \\
\midrule
DFR & 0.0147 & [-0.0249, 0.0656] & [-0.0145, 0.0884] & 0.750000 & Pass & Fail & Fail \\
Online logistic & -0.0434 & [-0.1054, 0.0000] & [-0.1387, 0.0000] & 0.250000 & Pass & Pass & Pass \\
\bottomrule
\end{tabular}
}
\par
}
\newcommand{\RunSixVerifiedPNNLCohortControlTable}{\noindent\resizebox{\linewidth}{!}{%
\begin{tabular}{llrrrrp{0.28\linewidth}}
\toprule
Plot ID & Cohort ID & Basis & d & r & Calibration pair & QASM/control status \\
\midrule
C01 & pgh\_d3\_r1\_x\_18\_11\_12\_13\_14 & X & 3 & 1 & edb355d8815d422dad137dca3e9c0a2b1034a64d1c58123b057c9cad8c667bc1--a5c3cd9a6a4a00f73eb0109e1690261ef335713605f06026b036f3758d29d694 & Raw and normalized QASM differ; circuit + hardware shift; no equal-QASM control \\
C02 & pgh\_d3\_r1\_x\_111\_98\_91\_92\_93 & X & 3 & 1 & edb355d8815d422dad137dca3e9c0a2b1034a64d1c58123b057c9cad8c667bc1--c20b55cb3b2ae3b5d1d16fbd540800fca9af9b2338f8c94035fff72db0e1081a & Raw and normalized QASM differ; circuit + hardware shift; no equal-QASM control \\
C03 & pgh\_d3\_r3\_x\_14\_13\_12\_11\_18 & X & 3 & 3 & c20b55cb3b2ae3b5d1d16fbd540800fca9af9b2338f8c94035fff72db0e1081a--a5c3cd9a6a4a00f73eb0109e1690261ef335713605f06026b036f3758d29d694 & Raw and normalized QASM differ; circuit + hardware shift; no equal-QASM control \\
C04 & pgh\_d3\_r3\_x\_115\_114\_113\_119\_133 & X & 3 & 3 & edb355d8815d422dad137dca3e9c0a2b1034a64d1c58123b057c9cad8c667bc1--c20b55cb3b2ae3b5d1d16fbd540800fca9af9b2338f8c94035fff72db0e1081a & Raw and normalized QASM differ; circuit + hardware shift; no equal-QASM control \\
C05 & pgh\_d3\_r3\_z\_151\_152\_153\_154\_155 & Z & 3 & 3 & 064bde2701d069ba9ccbe1c4b7d970796e901f09b32146d359fa02e2abe4e053--a5c3cd9a6a4a00f73eb0109e1690261ef335713605f06026b036f3758d29d694 & Raw and normalized QASM differ; circuit + hardware shift; no equal-QASM control \\
C06 & pgh\_d5\_r1\_x\_18\_11\_12\_13\_14\_15\_19\_35\_34 & X & 5 & 1 & edb355d8815d422dad137dca3e9c0a2b1034a64d1c58123b057c9cad8c667bc1--a5c3cd9a6a4a00f73eb0109e1690261ef335713605f06026b036f3758d29d694 & Raw and normalized QASM differ; circuit + hardware shift; no equal-QASM control \\
C07 & pgh\_d5\_r1\_z\_18\_11\_12\_13\_14\_15\_19\_35\_34 & Z & 5 & 1 & c20b55cb3b2ae3b5d1d16fbd540800fca9af9b2338f8c94035fff72db0e1081a--a5c3cd9a6a4a00f73eb0109e1690261ef335713605f06026b036f3758d29d694 & Raw and normalized QASM differ; circuit + hardware shift; no equal-QASM control \\
C08 & pgh\_d5\_r5\_x\_34\_35\_19\_15\_14\_13\_12\_11\_18 & X & 5 & 5 & edb355d8815d422dad137dca3e9c0a2b1034a64d1c58123b057c9cad8c667bc1--a5c3cd9a6a4a00f73eb0109e1690261ef335713605f06026b036f3758d29d694 & Raw and normalized QASM differ; circuit + hardware shift; no equal-QASM control \\
C09 & pgh\_d7\_r3\_z\_80\_81\_82\_83\_84\_85\_77\_65\_66\_67\_57\_47\_46 & Z & 7 & 3 & c20b55cb3b2ae3b5d1d16fbd540800fca9af9b2338f8c94035fff72db0e1081a--a5c3cd9a6a4a00f73eb0109e1690261ef335713605f06026b036f3758d29d694 & Raw and normalized QASM differ; circuit + hardware shift; no equal-QASM control \\
C10 & pgh\_d7\_r7\_z\_138\_151\_150\_149\_148\_147\_146\_145\_144\_143\_142\_141\_140 & Z & 7 & 7 & c20b55cb3b2ae3b5d1d16fbd540800fca9af9b2338f8c94035fff72db0e1081a--a5c3cd9a6a4a00f73eb0109e1690261ef335713605f06026b036f3758d29d694 & Raw and normalized QASM differ; circuit + hardware shift; no equal-QASM control \\
C11 & pgh\_d9\_r5\_x\_153\_154\_155\_139\_135\_134\_133\_119\_113\_114\_115\_99\_95\_94\_93\_92\_91 & X & 9 & 5 & b239f741f884c94069006f3043e2553bb70f45da4ecdb6623e41ffbbebcde5a7--c20b55cb3b2ae3b5d1d16fbd540800fca9af9b2338f8c94035fff72db0e1081a & Raw and normalized QASM differ; circuit + hardware shift; no equal-QASM control \\
\bottomrule
\end{tabular}
}
\par
}
\newcommand{\RunSixVerifiedPNNLRandomizationTable}{\noindent\textbf{Panel A: path--state alarm fractions.}\par
\noindent\resizebox{\linewidth}{!}{%
\begin{tabular}{llrrrrrr}
\toprule
Plot ID & Cohort ID & State & DFR & Online logistic & S-PACE sparse & S-PACE spectral & S-PACE composite \\
\midrule
C01 & pgh\_d3\_r1\_x\_18\_11\_12\_13\_14 & 0 & 0.000000 & 0.000000 & 0.000000 & 0.000000 & 0.000000 \\
C01 & pgh\_d3\_r1\_x\_18\_11\_12\_13\_14 & 1 & 0.000000 & 0.000000 & 0.000000 & 0.000000 & 0.000000 \\
C02 & pgh\_d3\_r1\_x\_111\_98\_91\_92\_93 & 0 & 0.003906 & 0.000000 & 0.000000 & 0.000000 & 0.000000 \\
C02 & pgh\_d3\_r1\_x\_111\_98\_91\_92\_93 & 1 & 0.000000 & 0.000000 & 0.007812 & 0.003906 & 0.007812 \\
C03 & pgh\_d3\_r3\_x\_14\_13\_12\_11\_18 & 0 & 0.000000 & 0.000000 & 0.003906 & 0.007812 & 0.007812 \\
C03 & pgh\_d3\_r3\_x\_14\_13\_12\_11\_18 & 1 & 0.000000 & 0.000000 & 0.003906 & 0.003906 & 0.007812 \\
C04 & pgh\_d3\_r3\_x\_115\_114\_113\_119\_133 & 0 & 0.000000 & 0.000000 & 0.007812 & 0.003906 & 0.003906 \\
C04 & pgh\_d3\_r3\_x\_115\_114\_113\_119\_133 & 1 & 0.003906 & 0.000000 & 0.007812 & 0.007812 & 0.011719 \\
C05 & pgh\_d3\_r3\_z\_151\_152\_153\_154\_155 & 0 & 0.000000 & 0.000000 & 0.000000 & 0.000000 & 0.000000 \\
C05 & pgh\_d3\_r3\_z\_151\_152\_153\_154\_155 & 1 & 0.000000 & 0.000000 & 0.007812 & 0.003906 & 0.003906 \\
C06 & pgh\_d5\_r1\_x\_18\_11\_12\_13\_14\_15\_19\_35\_34 & 0 & 0.000000 & 0.003906 & 0.015625 & 0.011719 & 0.015625 \\
C06 & pgh\_d5\_r1\_x\_18\_11\_12\_13\_14\_15\_19\_35\_34 & 1 & 0.000000 & 0.000000 & 0.000000 & 0.000000 & 0.000000 \\
C07 & pgh\_d5\_r1\_z\_18\_11\_12\_13\_14\_15\_19\_35\_34 & 0 & 0.000000 & 0.000000 & 0.003906 & 0.007812 & 0.007812 \\
C07 & pgh\_d5\_r1\_z\_18\_11\_12\_13\_14\_15\_19\_35\_34 & 1 & 0.000000 & 0.000000 & 0.000000 & 0.003906 & 0.003906 \\
C08 & pgh\_d5\_r5\_x\_34\_35\_19\_15\_14\_13\_12\_11\_18 & 0 & 0.003906 & 0.000000 & 0.000000 & 0.003906 & 0.003906 \\
C08 & pgh\_d5\_r5\_x\_34\_35\_19\_15\_14\_13\_12\_11\_18 & 1 & 0.000000 & 0.000000 & 0.003906 & 0.003906 & 0.003906 \\
C09 & pgh\_d7\_r3\_z\_80\_81\_82\_83\_84\_85\_77\_65\_66\_67\_57\_47\_46 & 0 & 0.000000 & 0.000000 & 0.000000 & 0.003906 & 0.003906 \\
C09 & pgh\_d7\_r3\_z\_80\_81\_82\_83\_84\_85\_77\_65\_66\_67\_57\_47\_46 & 1 & 0.000000 & 0.000000 & 0.003906 & 0.000000 & 0.000000 \\
C10 & pgh\_d7\_r7\_z\_138\_151\_150\_149\_148\_147\_146\_145\_144\_143\_142\_141\_140 & 0 & 0.000000 & 0.000000 & 0.000000 & 0.000000 & 0.000000 \\
C10 & pgh\_d7\_r7\_z\_138\_151\_150\_149\_148\_147\_146\_145\_144\_143\_142\_141\_140 & 1 & 0.000000 & 0.000000 & 0.000000 & 0.000000 & 0.000000 \\
C11 & pgh\_d9\_r5\_x\_153\_154\_155\_139\_135\_134\_133\_119\_113\_114\_115\_99\_95\_94\_93\_92\_91 & 0 & 0.000000 & 0.000000 & 0.003906 & 0.000000 & 0.000000 \\
C11 & pgh\_d9\_r5\_x\_153\_154\_155\_139\_135\_134\_133\_119\_113\_114\_115\_99\_95\_94\_93\_92\_91 & 1 & 0.000000 & 0.000000 & 0.007812 & 0.003906 & 0.000000 \\
\bottomrule
\end{tabular}
}
\par

\noindent\textbf{Panel B: method summaries over 256 replicates.}\par
\noindent\resizebox{\linewidth}{!}{%
\begin{tabular}{lrrr}
\toprule
Method & Episode alarm fraction & Alarmed episodes/replicate histogram & Maximum log-e (min/median/max) \\
\midrule
DFR & 0.000533 & 0:253, 1:3 & 0.6146 / 1.7825 / 6.2042 \\
Online logistic & 0.000178 & 0:255, 1:1 & 0.3544 / 1.2470 / 4.8024 \\
S-PACE sparse & 0.003551 & 0:237, 1:18, 2:1 & 0.8832 / 2.4226 / 8.9197 \\
S-PACE spectral & 0.003196 & 0:238, 1:18 & 0.8949 / 2.2424 / 6.9953 \\
S-PACE composite & 0.003729 & 0:236, 1:19, 2:1 & 0.8897 / 2.3429 / 8.2804 \\
\bottomrule
\end{tabular}
}
\par
}
\newcommand{\RunSixVerifiedResourceLedgerTable}{\noindent\resizebox{\linewidth}{!}{%
\begin{tabular}{lrrrrp{0.3\linewidth}}
\toprule
Stage & Paired-shot updates & Physical shots & Role updates & Detector bits & Timing/memory scope \\
\midrule
Google fit/warmup & 5,000 & 10,000 & 255,000 & 12,240,000 & Input accounting \\
Google threshold clone & 5,000 & 10,000 & 255,000 & 12,240,000 & Input accounting \\
Google held detector replay & 20,000 & 40,000 & 1,020,000 & 48,960,000 & Joint-pipeline runs 1312.527/1327.993/1339.677 s; median 1327.993 s; process-wide peak 1194.4 MiB \\
Google randomization surrogates & 1,280,000 & Resampled; no new exposure & 65,280,000 & No new exposure & 256 replicates; canonical merge only 8.956 s, 1 worker, process-wide peak 232.0 MiB; external shard concurrency not inferred \\
PNNL fit across 22 state streams & 8,184 & 16,368 & 24,552 & 212,784 & Whole-arm accounting; eigendecompositions 9,090 \\
PNNL surveillance & 24,556 & 49,112 & 73,668 & 638,400 & Whole-run 6900.189 s; held processing 6898.060 s; process-wide peak 1445.7 MiB; eigendecompositions 27,648 \\
PNNL threshold-bootstrap surrogates & 502,906,880 & Resampled; no new exposure & 1,508,720,640 & No new exposure & 4,096 replicates per path/state/method; eigendecompositions 226,492,416 \\
PNNL paired-swap surrogates & 20,953,600 & Resampled; no new exposure & 62,860,800 & No new exposure & 256 replicates; eigendecompositions 9,400,320 \\
\bottomrule
\end{tabular}
}
\par

\par\smallskip\noindent All three Google replay timings are for the canonical joint detector pipeline. The Google-randomization timing is for canonical merge only, not shard computation; external shard concurrency is not inferred. PNNL reports both whole-run and held-value processing wall time. Peak RSS is process-wide; surrogate rows are resampling work, not new physical-shot or detector-bit exposure. No per-method timing, memory, or relative-speed claim is reported.
}
\newcommand{\RunSixVerifiedGoogleEventFigure}{\noindent\includegraphics[width=\linewidth]{\RunSixGeneratedDir/google_event_alerts.pdf}}
\newcommand{\RunSixVerifiedGoogleRiskFigure}{\noindent\includegraphics[width=\linewidth]{\RunSixGeneratedDir/google_risk_coverage.pdf}}
\newcommand{\RunSixVerifiedGoogleRandomizationFigure}{\noindent\includegraphics[width=\linewidth]{\RunSixGeneratedDir/google_randomization_proper_prior.pdf}}
\newcommand{\RunSixVerifiedPNNLFigure}{\noindent\includegraphics[width=\linewidth]{\RunSixGeneratedDir/pnnl_delay_forest.pdf}}

\else
\fi
\newcommand{\RunSixClaimSentence}{%
  \ifrunresultbundle
    \RunSixVerifiedClaim%
  \else
    \RunSixPending{CLAIM_SENTENCE}%
  \fi
}

\newcommand{\E}{\mathbb E}
\newcommand{\Prb}{\mathbb P}
\newcommand{\KL}{\operatorname{KL}}
\newcommand{\Tr}{\operatorname{Tr}}
\newcommand{\rank}{\operatorname{rank}}
\newcommand{\cF}{\mathcal F}
\newcommand{\cC}{\mathcal C}
\newcommand{\one}{\mathbf 1}
\newcommand{\SPace}{S-PACE}
\newcommand{\DFR}{DFR}
\newcommand{\PNNL}{PNNL}

\newtheorem{theorem}{Theorem}
\newtheorem{proposition}[theorem]{Proposition}
\newtheorem{lemma}[theorem]{Lemma}
\theoremstyle{definition}
\newtheorem{definition}[theorem]{Definition}
\theoremstyle{remark}
\newtheorem{remark}[theorem]{Remark}

\setlength{\emergencystretch}{3em}
\renewcommand{\arraystretch}{1.12}

\begin{document}

\title{Predictable Sparse and Spectral Contrast Monitoring of Real QEC
Syndromes: A Predeclared Benchmark with a Disclosed Post-Detector Repair}

\author{Rongzhou (Lachlan) Chen}
\ifrunrevtex
  \affiliation{Independent research}
  \email{lachen@connect.hku.hk}
\else
  \date{Independent research\\Pre-results working paper, 28 July 2026}
\fi

\newcommand{\RunSixAbstractBody}{%
Quantum-error-correction (QEC) records can be nonstationary and burst-prone,
but detector firing rate is already a strong diagnostic and elaborate
feature maps need not improve operational decisions.  We introduce
\SPace{}, an initially preregistered composite of past-only sparse and spectral
contrast witnesses for classical syndrome streams.  Bounded predictable
scores are converted to betting factors and combined with fixed priors over
features, time scales, bet sizes, and 51 round-role experts.  One complete
paired shot---not one role---is the formal time step.  Exact false-alarm
validity is reserved for an independently randomized complete-pair design;
the natural Google replay and the constructed \PNNL{}/IBM snapshot
boundaries are evaluated empirically.  The locked benchmark compares
\SPace{} with paired detector firing rate, within-shot Page--CUSUM,
diagonal detector surprise, shrinkage Hotelling monitoring, and an online
logistic witness receiving the same feature bank.  Detector scores were
frozen under the original pre-access chain.  After that detector run, a
source-derived consumer-validation repair was separately ratified before
decoder outcomes, PNNL held payloads, or a completed randomization replicate
were available.  Detector numeric diagnostics had been exposed during
diagnosis but were not used to select the repair, and the detector was
neither modified nor rerun.  A conjunctive two-arm rule prevents a favorable
auxiliary analysis from rescuing a failed primary test.  \RunSixClaimSentence
}

\ifrunrevtex
\begin{abstract}
\RunSixAbstractBody
\end{abstract}
  \keywords{quantum error correction, sequential change detection,
  e-processes, preregistration, sparse features, spectral contrasts}
\maketitle
\else
\maketitle
\begin{abstract}
\RunSixAbstractBody
\end{abstract}
  \noindent\textbf{Keywords:} quantum error correction; sequential change
  detection; e-processes; preregistration; sparse features; spectral
  contrasts
\fi

\section{Introduction}
\label{sec:introduction}

Repeated syndrome measurements serve two distinct functions in QEC: they
support decoding, and they expose changes in the physical error process.
Detector rates, detector correlations, and occasional bursts are established
hardware diagnostics~\cite{google2023surface,mcewen2022bursts}.  Modern work
also adapts noise estimates and decoding graphs under drift
~\cite{wang2023dgr,bhardwaj2026drifting,poster2026reloqate}.  A new monitor
therefore earns its complexity only if it improves a fixed operational
endpoint at matched observations, calibration, and alert budget.

The motivating idea is to follow coordinates or operator directions with
large \emph{between-stream contrast}, rather than directions selected only
for pooled variance.  Sparse discriminative projections and tailored
high-dimensional change monitors already embody this broad principle
~\cite{tveten2019sparse}.  Likewise, predictable betting, fixed mixtures of
e-values, Page--CUSUM, Hotelling statistics, and the Jordan positive-part
construction are established foundations
~\cite{shafer2011testmartingales,vovk2021evalues,
waudbysmith2024betting,page1954cusum,hotelling1931ratio,
helstrom1969detection}.  The contribution pursued here is their
auditable integration into an initially preregistered real-QEC protocol,
followed by a disclosed post-detector, pre-outcome validator amendment:

\begin{enumerate}
  \item a fixed raw-check plus pair-correlation feature bank and a
  positive-semidefinite spectral representation;
  \item past-only sparse and spectral witnesses whose current scores remain
  bounded;
  \item a complete-shot mixture that does not count 51 correlated roles as
  51 fresh randomization revelations;
  \item separate proper-prior e-process and Shiryaev--Roberts (SR)
  accumulators, preserving the distinction between lifetime crossing
  probability and average run length (ARL);
  \item same-information rate, likelihood-surprise, covariance, CUSUM, and
  logistic controls; and
  \item a detector/outcome freeze and a falsifiable conjunctive empirical gate.
\end{enumerate}

Run 6 processes classical bits produced by fixed QEC measurements.  It is
not a quantum algorithm, a new measurement, or a claim of quantum speedup.
It cannot beat a correctly specified same-information likelihood ratio in
expected log growth, nor can it reveal physical changes that leave the
archived syndrome law unchanged.  These boundaries are part of the method,
not qualifications added after results.

\section{Observation Level, Related Work, and Evidence Boundary}
\label{sec:boundary}

\subsection{Three observation levels}

The experiment begins after the hardware measurement.  Table
~\ref{tab:levels} separates the three levels that must not be conflated.

\begin{table}[ht]
\centering
\caption{Observation levels and Run 6 access.}
\label{tab:levels}
\small
\begin{tabularx}{\linewidth}{@{}p{0.20\linewidth}X p{0.19\linewidth}@{}}
\toprule
Level & Object and relevant theory & Run 6 access\\
\midrule
Quantum acquisition & Unmeasured state/channel and measurement choice;
Helstrom or quantum quickest-change limits
~\cite{helstrom1969detection,fanizza2023quantumchange}. & No.\\
Classical stream & Fixed detector bits in archive order; likelihood,
CUSUM, covariance, and predictable betting. & Yes.\\
Retrospective utility & Observable flips and decoder predictions used for
risk--coverage after detector freeze. & Outcome join only.\\
\bottomrule
\end{tabularx}
\end{table}

Classical-shadow sequential detection requires randomized measurement
settings and a known inverse shadow channel
~\cite{huang2020shadows,zecchin2026escd}.  Those acquisition records are not
present in the fixed syndrome archives.  A numerical proxy would therefore
compare different experiments and is excluded.

\subsection{Sequential inference and change detection}

Nonnegative test martingales and e-values provide anytime evidence under
declared nulls~\cite{shafer2011testmartingales,vovk2021evalues,
howard2021confidence}.  Predictable betting gives concrete bounded-mean
constructions~\cite{waudbysmith2024betting}.  E-detectors provide CUSUM- and
SR-type ARL formulations~\cite{shin2024edetectors}; they do not turn every SR
threshold into a lifetime false-alarm probability.  Sequential
randomization e-values are also prior art~\cite{zampieri2025randomization}.
Run 6 applies these ideas with one orientation coin per complete QEC shot.

The empirical controls deliberately span simple and structured alternatives:
detector firing rate, Page's CUSUM~\cite{page1954cusum}, diagonal Bernoulli
surprise, Hotelling geometry with Ledoit--Wolf shrinkage
~\cite{hotelling1931ratio,ledoitwolf2004covariance}, and a pairwise online
logistic learner that receives the same 300 coordinates as the sparse
witness.  Published detector-likelihood work motivates, but is not
identical to, the diagonal surprise control used here
~\cite{hesner2025likelihood}.

\subsection{Evidence classes}

We use four labels:

\begin{enumerate}
  \item \textbf{Exact design audit}: complete observed pairs receive
  independently randomized orientation.
  \item \textbf{Natural approximate event}: archived Google order contains
  an author-identified, approximately located event.
  \item \textbf{Constructed boundary}: real-hardware \PNNL{}/IBM snapshot
  cohorts are concatenated by the benchmark.
  \item \textbf{Retrospective utility}: frozen detector ranks are joined to
  decoder mismatches.
\end{enumerate}

Only the first label supports the declared conditional randomization
guarantee.  The other arms supply empirical evidence at frozen operating
points.

\section{Preregistered Data Roles and Experimental Unit}
\label{sec:protocol}

\subsection{Google repetition-code replay}

The primary archive is the Google Quantum AI deposit accompanying its
surface-code study~\cite{google2022data,google2023surface}.  The locked
experiment
\texttt{repetition\_code\_bZ\_d25\_r50\_center\_5\_5} contains 500,000
shots, 50 syndrome rounds, and 1,224 declared detectors.  Detector
coordinates from the ideal Stim circuit define a permutation into 51 round
roles by 24 physical check coordinates.  The parser does not assume that a
raw reshape preserves the physical order.

The immutable split is shown in Table~\ref{tab:split}.  Held monitored shot
\(40000+q\) is paired one-to-one with reference shot \(20000+q\), for
\(q=0,\ldots,19999\).  Every complete pair consumes two archived physical
shots.  A role state is updated once per shot and is never shared with
another role.

\begin{table}[ht]
\centering
\caption{Frozen zero-based Google shot split.}
\label{tab:split}
\small
\begin{tabular}{@{}lll@{}}
\toprule
Indices & Role & Permitted use\\
\midrule
\([0,5000)\), \([10000,15000)\) & paired fit & model state\\
\([5000,10000)\), \([15000,20000)\) & threshold & operating points\\
\([20000,40000)\) & held reference & one per test shot\\
\([40000,60000)\) & held monitor & primary replay\\
\([60000,500000)\) & future & untouched in primary\\
\bottomrule
\end{tabular}
\end{table}

The primary event window is \([57750,57800)\), with frozen narrow and wide
sensitivity windows \([57770,57780)\) and \([57725,57825)\).  The archive
documentation localizes a large event only approximately; no exact physical
onset or wall-clock delay is inferred.

\subsection{\PNNL{}/IBM auxiliary arm}

The auxiliary data are the versioned \PNNL{} release of IBM
repetition-code experiments~\cite{stein2026film,stein2026dataset}.  Physical
paths are reconstructed from each state-specific QASM.  The independent
summary unit is a path/snapshot, not an individual syndrome bit.  Every
comparison is called a \emph{constructed boundary between real-hardware
snapshot cohorts}.  It is not interpreted as a natural time series; changed
QASM is reported as circuit-and-hardware domain shift.

\subsection{One randomization unit, one formal time}

For complete paired shot \(u\), a single random bit swaps its A/B
orientation.  The same bit acts on all roles \(r=0,\ldots,50\).  All role
scores use pre-shot states, after which the 51 role-component experts are
updated in parallel.  Formal time is \(u\), not \(51u+r\).  The flattened
cycle index is retained only for empirical alert localization.

\begin{figure}[ht]
\centering
\fbox{\begin{minipage}{0.92\linewidth}
\centering
\textbf{Frozen causal sequence}\\[3pt]
fit pairs \(\rightarrow\) checkpoint clone
\(\rightarrow\) disjoint threshold replay
\(\rightarrow\) held detector replay
\(\rightarrow\) hash detector artifacts
\(\rightarrow\) outcome join\\[4pt]
One complete pair \(\rightarrow\) 51 simultaneous role scores
\(\rightarrow\) one formal update.
\end{minipage}}
\caption{Protocol timeline.  Threshold-state updates do not flow into held
monitoring, and outcome labels cannot affect detector scores or ties.}
\label{fig:timeline}
\end{figure}

\section{\SPace{} Construction}
\label{sec:method}

\subsection{Bounded global, sparse, and spectral representations}

For one cycle let \(e\in\{0,1\}^{24}\), \(z=\one-2e\), and
\[
  g(e)=\frac1{24}\sum_{j=1}^{24}e_j.
\]
The fixed sparse feature bank is
\[
  \phi(e)=
  \left(e_1,\ldots,e_{24},
  \left\{\frac{1+z_i z_j}{2}\right\}_{1\le i<j\le24}\right)
  \in[0,1]^{300}.
\tag{1}
\]
All 276 check pairs are included, preventing an event-dependent adjacency
choice.  The spectral representation is
\[
  R(e)=\frac{zz^{\mathsf T}}{24},\qquad
  R(e)\succeq0,\quad \Tr R(e)=1.
\tag{2}
\]
For paired reference/monitor cycles,
\[
  D^\phi_{u,r}=\phi(e^B_{u,r})-\phi(e^A_{u,r}),\qquad
  \Delta_{u,r}=R(e^B_{u,r})-R(e^A_{u,r}).
\tag{3}
\]

The separately designated M0 control uses
\(g(e^B)-g(e^A)\).  It is not a component of the fixed proposed method
\(S\), which is exactly the M4/M5 composite defined below.  Reporting M0
alongside \(S\) supplies a directly interpretable global-rate baseline; a
pure orbit-centered relative representation would be pathwise blind to an
additive common mode, as Sec.~\ref{sec:guarantees} makes explicit.

\subsection{Sparse analytical witness}

For each role, sparse witnesses retain exponentially weighted, past-only
contrasts with half-lives \(h\in\{4,16,64,256\}\).  For
\(k\in\{1,4,16,64\}\), the witness assigns weight \(1/k\) to the signed
copies of the \(k\) largest absolute historical contrasts, using a fixed
magnitude/index tie rule.  Its \(\ell_1\) norm is at most one, so the current
score \(s^{(4)}_{u,r,h,k}=w_{u-1,r,h,k}^{\mathsf T}D^\phi_{u,r}\) lies in
\([-1,1]\).  Scoring precedes the EWMA update.

\subsection{Spectral analytical-observable-contrast witness}

For \(h\in\{4,16,64\}\), let \(\overline\Delta_{u-1,r,h}\) be the past-only
EWMA operator contrast.  The stored matrix is symmetrized in float64.
Effects are recomputed at prior role-update counts \(8,16,24,\ldots\) and
first used on the following observation, with absolute eigentolerance
\(\tau_{\mathrm{eig}}=10^{-10}\).  At a locked refresh, the full effect is
\[
  E^{+}_{u-1,r,h}
  =\Pi_+(\overline\Delta_{u-1,r,h};\tau_{\mathrm{eig}}),
\tag{4}
\]
where
\(\Pi_+(H;\tau)=\sum_{\lambda_i(H)>\tau}v_iv_i^{\mathsf T}\), or zero if
the index set is empty.  For the
rank-one variant, if \(\lambda_{\max}\le\tau_{\mathrm{eig}}\) the effect is
zero.  Otherwise form the projector \(P_{\max}\) of all eigenvalues within
\(\tau_{\mathrm{eig}}\) of \(\lambda_{\max}\); choose the smallest coordinate
\(j\) whose diagonal entry is within tolerance of
\(\max_i(P_{\max})_{ii}\); and set
\[
  v=\frac{P_{\max}e_j}{\|P_{\max}e_j\|_2},\qquad E^{(1)}=vv^{\mathsf T}.
\]
The first entry of \(v\) with magnitude above tolerance is made positive,
although its projector is sign invariant.  Both effects obey
\(0\preceq E\preceq I\);
therefore
\[
  s^{(5)}_{u,r,h,\mathrm{rank}}
  =\Tr(E_{u-1,r,h,\mathrm{rank}}\Delta_{u,r})\in[-1,1].
\tag{5}
\]
The positive-part variational identity is the standard Jordan/Helstrom
calculation, applied here to an already observed classical feature matrix;
it is not a physical Helstrom measurement.

\subsection{Factors, fixed priors, and empirical score}

Every bounded predictable score is converted to
\[
  L_{u,r,c}=1+\beta_c s_{u,r,c},
  \qquad \beta_c\in\{0.1,0.3,0.6,0.9\}.
\tag{6}
\]
M4 has 64 uniform components over \((h,k,\beta)\); M5 has 24 over
\((h,\mathrm{rank},\beta)\).  The exact \SPace{} composite assigns one half
of its prior to each branch and remains a fixed mixture.  It never selects
the better branch from held performance.

The empirically calibrated cycle statistic is separately fixed as
\[
  Z^{(S)}_{q,r}
  =\max\{Z^{(4)}_{q,r},Z^{(5)}_{q,r}\},
\tag{7}
\]
where each branch maximum ranges only over its predeclared components.  This
maximum is \emph{not} called an e-factor; its full multiplicity is calibrated
on the disjoint threshold block.

\section{Guarantees and Impossibility Boundaries}
\label{sec:guarantees}

\subsection{Predictable factors}

\begin{proposition}[bounded predictable factor]
\label{prop:factor}
Let \(s_u\in[-1,1]\) satisfy
\(\E_0[s_u\mid\cF_{u-1}]=0\), and let
\(\beta_u\in[-1,1]\) be \(\cF_{u-1}\)-measurable.  Then
\(L_u=1+\beta_u s_u\) is nonnegative and
\(\E_0[L_u\mid\cF_{u-1}]=1\).
\end{proposition}

\begin{proof}
The range assumptions give \(L_u\ge0\), and conditional linearity gives
\(\E_0[L_u\mid\cF_{u-1}]=1+\beta_u\E_0[s_u\mid\cF_{u-1}]=1\).
\end{proof}

This proposition is a direct predictable-betting specialization, not a new
martingale theorem.  It applies to the independently randomized orientation
because an antisymmetric current-pair score is conditionally centered given
the unordered pair and past state.  Fixed natural hardware orientation is
not made exchangeable by this notation.

\subsection{Complete-shot expert mixture}

\begin{proposition}[shared-orientation roles]
\label{prop:roles}
Suppose one orientation coin is drawn for complete paired shot \(u\), and
every role-component score is computed from its pre-shot state.  If each
\((r,c)\) construction is a conditional e-factor with respect to the
pre-coin filtration, then every fixed convex mixture of the corresponding
role-component wealth processes is an e-process.  Serially treating the 51
roles as fresh coin revelations is not justified.
\end{proposition}

\begin{proof}
Each expert wealth is nonnegative with conditional expectation no larger
than its previous value.  Conditional linearity preserves this property
under fixed convex weights despite arbitrary dependence between experts.
After the first role reveals the shared orientation, the remaining roles do
not receive a fresh randomization bit, so a 51-step product needs an
additional conditional argument that the design does not supply.
\end{proof}

For horizon \(H\), the locked proper-prior recursion is
\[
  A_{0,r,c}=0,\qquad
  A_{u,r,c}=L_{u,r,c}\!\left(A_{u-1,r,c}+\frac1H\right),
\tag{8}
\]
\[
  E_u=\frac{H-u}{H}
  +\frac1{51}\sum_{r=0}^{50}\sum_c\pi_c A_{u,r,c}.
\tag{9}
\]
Here \(H=20{,}000\) for held replay and \(H=5{,}000\) for the randomization
threshold replay.

\subsection{Common-mode blind spot and global diagnostic}

\begin{proposition}[relative-projection blind spot]
\label{prop:common}
Let \(P_{\mathrm{rel}}=I-P_{\mathrm{triv}}\) remove the invariant spatial
mode.  If a change is exactly
\(\Phi'_u=\Phi_u+\one\otimes c_u\), with no clipping, nonlinear remapping, or
relative change, then
\[
  P_{\mathrm{rel}}\Phi'_u=P_{\mathrm{rel}}\Phi_u.
\]
Every detector measurable only with respect to the relative field has the
same path law before and after this change.
\end{proposition}

\begin{proof}
The additive term belongs to the invariant subspace and is annihilated by
\(P_{\mathrm{rel}}\).  Equality is pathwise, so equality of every
relative-field functional follows.
\end{proof}

This projection identity motivates reporting M0 as a separate, directly
interpretable global-rate diagnostic beside the fixed M4/M5 composite \(S\).
M0 is a control, not a claimed repair or an additional branch of \(S\), and
the identity is not a new statement about topological order.

\subsection{Closed-form maximum-difference coordinates}

\begin{proposition}[top-\(k\) capped contrast]
\label{prop:topk}
For mean contrast \(\delta\in\mathbb R^p\) and
\(k\in\{1,\ldots,p\}\), let
\(v=(\delta,-\delta)\) and
\[
  \cC_k=\left\{q\in\mathbb R^{2p}_+:
  \sum_{i=1}^{2p}q_i=1,\ q_i\le\frac1k\right\}.
\]
Then
\[
  \max_{q\in\cC_k}q^{\mathsf T}v
  =\frac1k\sum_{j=1}^{k}|\delta|_{(j)}.
\tag{10}
\]
\end{proposition}

\begin{proof}
A linear functional over the capped simplex assigns the largest allowed
mass to its largest coordinates.  The \(k\) largest coordinates of
\((\delta,-\delta)\) are the correctly signed copies of the \(k\) largest
absolute entries of \(\delta\).
\end{proof}

Equation~(10) solves the declared linear-gap objective.  It does not prove
minimum detection delay or universal feature-selection optimality.

\subsection{Spectral gap and likelihood ceiling}

\begin{proposition}[positive-part operator gap]
\label{prop:jordan}
For Hermitian \(\overline\Delta\),
\[
  \max_{0\preceq E\preceq I}\Tr(E\overline\Delta)
  =\Tr(\overline\Delta_+),
\]
attained by the unambiguous spectral projector
\(E=\Pi_+(\overline\Delta;0)\) onto the strictly positive eigenspace.
\end{proposition}

This standard variational identity maximizes a one-step linear gap, not
sequential log wealth.  Moreover, if \(P_1\ll P_0\) is a correctly specified
same-information alternative, every one-step e-factor \(L\) obeys
\[
  \E_1\log L\le\KL(P_1\Vert P_0),
\tag{11}
\]
with equality for the likelihood ratio \(dP_1/dP_0\).
Indeed,
\[
\E_1\log L
=\KL(P_1\Vert P_0)
+\E_1\log\!\left(L\frac{dP_0}{dP_1}\right)
\le\KL(P_1\Vert P_0)
\]
by Jensen and \(\E_0L\le1\).  Thus potential benefits of \SPace{} concern
robustness, transparent calibration, sparse interpretation, or model-fitting
cost---not violation of the likelihood ceiling
~\cite{neymanpearson1933tests}.

\subsection{Lifetime control versus average run length}

\begin{theorem}[declared sequential meanings]
\label{thm:sequential}
Under the exact conditional null, Eq.~(9) is a nonnegative
supermartingale with \(E_0=1\), so
\[
  \Prb_0\!\left(\sup_{u\le H}E_u\ge1/\alpha\right)\le\alpha.
\tag{12}
\]
Separately, if
\[
  R_{u,r,c}=(R_{u-1,r,c}+1)L_{u,r,c},\qquad
  R_u=\frac1{51}\sum_{r,c}\pi_cR_{u,r,c},
\tag{13}
\]
then the SR stopping rule at threshold \(\gamma\) has null ARL at least
\(\gamma\) \emph{complete paired shots}, subject to the standard
optional-sampling conditions.  A role or flattened cycle index can localize
an empirical alert within a shot but is not an ARL time unit.
\end{theorem}

These statements are not interchangeable.  When \(L\equiv1\), the SR
statistic advances like a clock and crosses a finite threshold.  Run 6 uses
\(\alpha=0.01\) for the reported proper-prior audit.  Equation~(13) records
the distinct SR/ARL construction and its formal unit; this manuscript does
not report a separate Google SR result.

\subsection{Accessible-data no-go}

\begin{theorem}[pushforward no-go]
\label{thm:nogo}
Let \(Y\) denote the archived syndrome path and let \(T(Y,U)\) be any
detector, stopping rule, or randomized post-processing, where auxiliary
randomness \(U\) has the same law under both hypotheses.  If the two physical
regimes induce the same law for \(Y\), they induce the same law for
\(T(Y,U)\).
\end{theorem}

\begin{proof}
The result is the data-processing identity for the pushforward of the common
law of \((Y,U)\).
\end{proof}

Theorem~\ref{thm:nogo} does not discover Wilson loops, toric-code logical
structure, or a string-theory result.  Wilson observables and toric-code
strings are established physical constructions
~\cite{wilson1974confinement,kitaev2003anyons}.  They are not measured in
the primary repetition-code archive.

\section{Locked Comparators and Evaluation}
\label{sec:evaluation}

\subsection{Methods and information parity}

Table~\ref{tab:methods} records the frozen method set.  M0C, M1, and M2 are
mandatory contextual controls but cannot rescue failure of the
predesignated M0/M3 comparisons.

\begin{table}[ht]
\centering
\caption{Frozen methods.  ``Exact'' refers only to the randomized complete-
pair design with the displayed factor construction.}
\label{tab:methods}
\small
\begin{tabularx}{\linewidth}{@{}lXl@{}}
\toprule
ID & Score or fitted object & Status\\
\midrule
M0 & paired detector firing-rate difference & exact-factor/empirical\\
M0C & within-shot two-sided Page--CUSUM & empirical\\
M1 & fixed diagonal Bernoulli surprise & exact-factor/empirical\\
M2 & role-centered shrinkage Hotelling score & empirical\\
M3 & past-only pairwise logistic, same \(\phi\) & exact-factor/empirical\\
M4 & sparse signed top-\(k\) witness & exact-factor/empirical\\
M5 & positive-eigenspace/rank-one witness & exact-factor/empirical\\
S & fixed M4/M5 prior; empirical max score & as declared above\\
\bottomrule
\end{tabularx}
\end{table}

\paragraph{M0 and M0C.}
With \(d^g_{q,r}=g(e^B_{q,r})-g(e^A_{q,r})\), M0 uses
\(s^{(0)}_{q,r}=d^g_{q,r}\), eight signed bets
\(\beta\in\{\pm0.1,\pm0.3,\pm0.6,\pm0.9\}\), and empirical score
\(Z^{(0)}_{q,r}=|s^{(0)}_{q,r}|\).  At shot \(q\), define the M0C channels
\(x_{r,0}=d^g_{q,r}\) and
\(x_{r,j}=e^B_{q,r,j}-e^A_{q,r,j}\) for raw checks
\(j\in\{1,\ldots,24\}\).  M0C resets at each shot and, for drifts
\(\kappa\in\{0.01,0.05,0.1\}\), applies
\[
 C^+_{\kappa,c}\leftarrow
 \max\{0,C^+_{\kappa,c}+x_{r,c}-\kappa\},\qquad
 C^-_{\kappa,c}\leftarrow
 \max\{0,C^-_{\kappa,c}-x_{r,c}-\kappa\}.
\]
Its score is \(\max_{\kappa,c}\{C^+_{\kappa,c},C^-_{\kappa,c}\}\);
it is a within-shot empirical scan, not an e-factor or cross-shot
quickest-change oracle.

\paragraph{M1 and M2.}
M1 freezes role/check probabilities from the fit pairs using Jeffreys
smoothing and clipping,
\[
 \widehat p_{rj}=\operatorname{clip}
 \left(\frac{c_{rj}+1/2}{10001},10^{-4},1-10^{-4}\right),
\]
forms the mean diagonal Bernoulli negative log likelihood \(\ell_r(e)\),
and scores
\[
 s^{(1)}_{q,r}=
 \frac{\ell_r(e^B_{q,r})-\ell_r(e^A_{q,r})}
 {\log\{(1-10^{-4})/10^{-4}\}}\in[-1,1].
\]
M2 uses the exact PCG64(610601) role-stratified 20,000-observation fit
subsample, role means \(\widehat\mu_r\), and one pooled symmetrized
Ledoit--Wolf precision \(\widehat\Omega\):
\[
 Z^{(2)}_{q,r}=\max\!\left\{0,
 (D^\phi_{q,r}-\widehat\mu_r)^{\mathsf T}\widehat\Omega
 (D^\phi_{q,r}-\widehat\mu_r)\right\}.
\]
M2 is empirical and never inherits an e-process guarantee.

\paragraph{M3 same-feature control.}
For \(\eta\in\{0.001,0.01,0.1\}\), M3 maintains a distinct
\(w_{r,\eta}\in\mathbb R^{300}\), no intercept, and L2 coefficient
\(\lambda=10^{-4}\).  From the pre-update
\(a=w_{r,\eta}^{\mathsf T}D^\phi_{q,r}\) and
\(2c=\max\{1,\|w_{r,\eta}\|_1\}\), it scores
\[
 s^{(3,\eta)}_{q,r}=\tanh(a/(2c))
\]
and only then applies the overflow-safe pairwise-logistic update
\[
 w_{r,\eta}\leftarrow w_{r,\eta}
 +\eta\,\sigma(-a)D^\phi_{q,r}
 -\eta\lambda w_{r,\eta}.
\]
The empirical maximum over the three frozen experts is threshold-calibrated;
the exact branch instead uses its fixed 12-component \((\eta,\beta)\)
mixture.

All methods receive the same raw detector records.  The 5,000-pair fit block
is common.  The disjoint 5,000-pair threshold clone selects the smallest
observed-or-infinite strict threshold whose replay emits at most two
notifications among 255,000 cycle-role opportunities.  At most one
notification is counted per shot; states continue updating through
within-shot suppression.  The complete threshold frontier is frozen before
held scoring.

\subsection{Primary event and risk--coverage endpoints}

The natural-event endpoints are: alarm inside each frozen event window,
first alarm shot and role, pre-event alerts, and event score rank.  Because
the event is a single cluster with approximate onset, cycle-level population
confidence claims are prohibited.

Only after detector artifacts are hashed are actual observable flips and
decoder predictions opened.  At the primary triage budget, each method ranks
the 20,000 monitor shots by frozen score, breaking ties by ascending archive
index.  The primary label is actual observable flip XOR correlated-matching
prediction; PyMatching is secondary.  The top 20 shots are selected without
labels.  This is retrospective veto/triage, not a new decoder.

\subsection{Post-detector repair chronology and dual provenance}
\label{sec:repair}

The detector implementation, manifest, and original ratification were
committed and pushed before detector-value access.  The detector replay then
read, scored, and serialized the held detector stream.  The first
randomization launch stopped in artifact-schema validation: all 32 planned
shards produced the same failure, no shard manifest and no completed
replicate.  Diagnosis exposed detector-derived numeric diagnostics.

The defect was a producer--consumer validation mismatch: the producer omitted
one redundant SR metadata key, while the consumers required it, and the
consumer's expected prior vector did not reproduce the producer's exact
normalization path.  The repair was fixed from source and schema evidence.
No detector-derived number selected the repair; no threshold, score, state,
seed, window, endpoint, or claim gate changed.  The original 231-artifact
detector bundle was reused byte-for-byte and was not rerun.

A second, post-detector and pre-outcome
implementation--manifest--ratification chain binds the exact allowlisted
consumer and verification diff, the original ratification, detector
manifest and artifact registry, 32 preserved failed attempts, environment,
thread settings, package lock, and runtime module origins.  Its access record
truthfully sets detector access and numeric diagnostic exposure to true,
repair selection from those numbers to false, decoder-outcome and PNNL held
access to false, and completed randomization to zero.  Randomization, PNNL,
the outcome join, and the final decision must bind this repair ratification;
the immutable detector manifest continues to bind only the original
ratification.

Thus the methods and gates were initially preregistered, but the end-to-end
execution is not described as detector-blind repair or pristine full
preregistration.  This chronology is a limitation and part of the
machine-checked publication contract.

\subsection{Locked Google uncertainty}

Risk--coverage stability uses a paired circular moving-block bootstrap of
complete monitored shots: block length 128, 2,000 replicates, and replicate
seed \(611000+b\).  Each resample draws 157 circular block starts, truncates
to 20,000 shots, resamples score/outcome rows together, and recomputes the
fixed top-\(B\) metrics and differences.  A separate descriptive threshold
bootstrap uses complete threshold shots, block length 128, 2,000 replicates,
40 block starts, and seeds \(613000+b\).  Neither bootstrap replaces the
deterministic threshold or changes the primary Boolean.  With one natural
event cluster, neither interval is called a population-level event-power
confidence interval.

\subsection{Exact randomization audit}

The threshold block is replayed under 256 predeclared random-swap seeds.  One
orientation bit applies to all 51 roles in a complete pair, all states reset
between replicates, and factors are mixed only at complete-shot time.  This
arm audits implementation and design-based calibration.  It neither
establishes natural hardware stationarity nor estimates event power.

\subsection{\PNNL{} path-level uncertainty}

Each of the 22 path--logical-state streams calibrates its empirical threshold
with 4,096 circular moving-block bootstrap replicates of \emph{complete
paired shots}, block length 32, and its frozen cohort/state/method seed.
This within-stream threshold bootstrap is not the cross-cohort uncertainty
unit.

For the auxiliary comparison, logical states 0 and 1 are first averaged
within each of the 11 physical path/snapshot pairs.  The primary uncertainty
analysis then bootstraps these 11 paired path effects 10,000 times with
PCG64 seed 612500.  A separate 10,000-draw sensitivity bootstrap clusters by
the five ordered calibration-hash pairs with seed 612501.  The paired
randomization analysis exhaustively enumerates all \(2^{11}\) sign flips of
the 11 path effects.  No bootstrap or randomization treats individual
syndrome bits as independent.

The exact swap audit is a different operation: seeds 610700 through 610955
create one A/B swap bit per complete paired shot, shared across all methods
and round roles; every path--state replicate resets fully, scores before
updating, and performs one role-expert mixture update per shot.  It validates
the design, not natural IBM drift.  QASM equality and inequality remain
visible so circuit-and-hardware domain shift cannot be mislabeled as
hardware-only drift.  The publication consumer validates all 110 ordered
path--state--method JSON rows and aggregate-cross-checks their alarm totals
and global maxima against hash-bound \(256\times5\) count and maximum-log-\(e\)
arrays.  Those aggregate arrays do not reconstruct the allocation of alarms
among individual path rows, so the audit is not described as an independent
row-level recomputation.

\subsection{Locked conjunctive empirical gate}

Let \(G\) be \texttt{google\_primary\_pass}.  It is true only if the fixed
composite:

\begin{enumerate}
  \item alarms inside \([57750,57800)\) at the validation-locked primary
  threshold;
  \item has at most nine pre-event alerts;
  \item was frozen with its resource ledger before outcome access;
  \item captures at least one more top-20 correlated-decoder mismatch than
  paired \DFR{};
  \item captures at least one more than same-feature online logistic;
  \item reports every frozen sensitivity and uncertainty summary; and
  \item receives no extra records or outcome labels.
\end{enumerate}

Let \(P\) be the single aggregate
\texttt{pnnl\_retention\_pass}.  Against \emph{each} of DFR and
same-feature online logistic, the fixed composite must have no more
pre-boundary false alarms over all 22 state streams and a strictly lower
equal-weight macro restricted-delay fraction after logical-state averaging.
Ties fail.  This is not a count or selection of ``retained paths''; all 11
path rows and both logical-state summaries remain in the report.  The
overall statement is the conjunction
\[
  C_{\mathrm{Run\,6}}=G\land P.
\tag{14}
\]
If either gate fails, the required conclusion is exactly ``no demonstrated
\SPace{} algorithmic advantage.''  A tie, a sparse heat map, or a successful
randomization audit cannot override Eq.~(14).

\subsection{Physical and computational resource ledger}

Observation exposure is reported for every method: calibration/reference
and monitor shots, QEC cycles, detector bits, feature dimension, and update
frequency.  Offline pairing consumes two physical archived shots per formal
update.  The implementation executes all methods in one detector runner, so
wall time and maximum resident set size are identifiable only for the
\emph{joint pipeline/process}.  They are not divided into unsupported
per-method timing or RSS claims.  The process-wide ledger additionally
records output bytes and high-level operation counts where the manifest
supports them.  Formal e-process/ARL time is always a complete paired shot;
role and cycle indices are only within-shot localization coordinates.

\section{Locked Results and Dual Provenance}
\label{sec:results}

\ifrunresultbundle\else
\noindent\textcolor{pending}{\textbf{Pre-results status.}}
Every red \texttt{TBD(KEY)} marker is a conditional fallback.  Final mode
accepts only fragments generated from the original freeze, repair chain, and
four completed hash-bound result manifests by the fail-closed publication
extractor; editing this manuscript cannot resolve a marker.
\fi

\subsection{Provenance and detector freeze}

The publication bundle verifies two distinct provenance lines.  The immutable
detector manifest binds the original pre-access ratification.  The
randomization, PNNL, outcome, and decision records additionally bind the
post-detector repair ratification and its truthful access record.  The
extractor rejects any attempt to describe the repair as detector-blind, any
detector rerun, or any completed randomization or held-outcome/PNNL access
before repair.  It verifies every derived artifact digest and recomputes both
empirical gates before emitting the claim sentence.  The bundle contract is
\ifrunresultbundle
\texttt{generated/publication\_bundle\_manifest.json}.
\else
\RunSixPending{PUBLICATION_BUNDLE_MANIFEST}.
\fi
Detector-only artifacts precede both the repair ratification and the
separately hash-authorized outcome join.

\subsection{Threshold frontier and locked uncertainty}

\begin{table}[ht]
\centering
\caption{Validation-locked thresholds and complete alert frontier.  Strict
crossing and one-notification-per-shot semantics remain fixed.}
\label{tab:frontier}
\small
\ifrunresultbundle
  \resizebox{\linewidth}{!}{\RunSixVerifiedGoogleThresholdFrontierTable}
\else
  \RunSixPending{GOOGLE_THRESHOLD_FRONTIER_TABLE}
\fi
\end{table}

\begin{table}[ht]
\centering
\caption{Descriptive Google threshold-bootstrap audit for all methods.
Complete-shot circular blocks have length 128; all 2,000 replicates and
seeds 613000--614999 are validated.  The summaries do not replace a frozen
threshold, enter a locked Boolean, or estimate natural-event power.}
\label{tab:google-threshold-bootstrap}
\small
\ifrunresultbundle
  \resizebox{\linewidth}{!}{\RunSixVerifiedGoogleThresholdBootstrapTable}
\else
  \RunSixPending{GOOGLE_THRESHOLD_BOOTSTRAP_TABLE}
\fi
\end{table}

\begin{table}[ht]
\centering
\caption{Locked Google S-PACE-minus-comparator uncertainty for both labels,
both comparators, budgets \(2,20,200\), and all four metrics, including valid
replicate counts.  Resampling uses complete shots/blocks and does not change
the deterministic primary Boolean; only primary-label budget-20 capture is
gate-relevant.}
\label{tab:google-uncertainty}
\small
\ifrunresultbundle
  \resizebox{\linewidth}{!}{\RunSixVerifiedGoogleUncertaintyTable}
\else
  \RunSixPending{GOOGLE_UNCERTAINTY_TABLE}
\fi
\end{table}

\subsection{Google natural-event replay}

\begin{table}[ht]
\centering
\caption{Frozen Google event results for all predeclared methods, including
the mandatory CUSUM, diagonal-surprise, and Hotelling controls.}
\label{tab:google}
\small
\ifrunresultbundle
  \resizebox{\linewidth}{!}{\RunSixVerifiedGoogleEventTable}
\else
  \RunSixPending{GOOGLE_EVENT_TABLE}
\fi
\end{table}

\begin{figure}[ht]
\centering
\ifrunresultbundle
  \RunSixVerifiedGoogleEventFigure
\else
  \fbox{\parbox[c][32mm][c]{0.90\linewidth}{\centering
  \RunSixPending{GOOGLE_EVENT_FIGURE}}}
\fi
\caption{Google natural-event causal replay.  The author-provided event
location is approximate; the plot does not define an exact onset.}
\label{fig:google}
\end{figure}

\subsection{Frozen decoder-outcome join}

\begin{table}[ht]
\centering
\caption{Risk--coverage results from frozen rankings for both decoder labels
and all locked budgets \(2,20,200\), with all eight methods, captured/total,
recall, precision, retained mismatch rate, coverage, intervals, and valid
replicate counts visible.  Only correlated-matching capture at budget 20
enters the primary gate.}
\label{tab:risk-budget}
\small
\ifrunresultbundle
  \resizebox{\linewidth}{!}{\RunSixVerifiedGoogleRiskBudgetTable}
\else
  \RunSixPending{GOOGLE_RISK_BUDGET_TABLE}
\fi
\end{table}

\begin{figure}[ht]
\centering
\ifrunresultbundle
  \RunSixVerifiedGoogleRiskFigure
\else
  \fbox{\parbox[c][32mm][c]{0.90\linewidth}{\centering
  \RunSixPending{GOOGLE_RISK_COVERAGE_FIGURE}}}
\fi
\caption{Retrospective utility of detector-only frozen rankings.}
\label{fig:risk}
\end{figure}

\subsection{Randomization-design calibration}

\begin{table}[ht]
\centering
\caption{Complete-shot proper-prior randomization audit.  No SR result is
inferred or reported.}
\label{tab:randomization}
\small
\ifrunresultbundle
  \resizebox{\linewidth}{!}{\RunSixVerifiedGoogleRandomizationTable}
\else
  \RunSixPending{GOOGLE_RANDOMIZATION_TABLE}
\fi
\end{table}

\begin{figure}[ht]
\centering
\ifrunresultbundle
  \RunSixVerifiedGoogleRandomizationFigure
\else
  \fbox{\parbox[c][32mm][c]{0.90\linewidth}{\centering
  \RunSixPending{GOOGLE_RANDOMIZATION_FIGURE}}}
\fi
\caption{Exact-design randomization audit.}
\label{fig:randomization}
\end{figure}

\subsection{\PNNL{}/IBM constructed-boundary arm}

\begin{table}[ht]
\centering
\caption{\PNNL{} exact paired-swap implementation audit.  The table covers
all 22 path--state episodes, five methods, and 256 seeds.  It is not a
natural-hardware null, enters neither empirical gate nor the overall
conjunction, and cannot override a failed empirical gate.}
\label{tab:pnnl-randomization}
\small
\ifrunresultbundle
  \RunSixVerifiedPNNLRandomizationTable
\else
  \RunSixPending{PNNL_RANDOMIZATION_AUDIT_TABLE}
\fi
\end{table}

\begin{table}[ht]
\centering
\caption{\PNNL{} equal-weight macro summaries.  Pre-boundary false-alarm
and miss counts cover all 22 path--state streams; no path is selected or
discarded by the retention decision.}
\label{tab:pnnl-macro}
\small
\ifrunresultbundle
  \resizebox{\linewidth}{!}{\RunSixVerifiedPNNLMacroTable}
\else
  \RunSixPending{PNNL_MACRO_PREFA_MISS_TABLE}
\fi
\end{table}

\begin{table}[ht]
\centering
\caption{\PNNL{} state-resolved results.  Each cell is
pre-boundary false-alarm indicator / miss indicator / restricted post-boundary
delay fraction.  All 22 ordered path--state rows and all five methods are
reported; no path or logical state is selected.}
\label{tab:pnnl-state-results}
\small
\ifrunresultbundle
  \resizebox{\linewidth}{!}{\RunSixVerifiedPNNLStateResultsTable}
\else
  \RunSixPending{PNNL_STATE_RESULTS_TABLE}
\fi
\end{table}

\begin{table}[ht]
\centering
\caption{\PNNL{} state-averaged cohort results.  Each cell is pre-false-alarm
mean / miss mean / restricted post-boundary delay fraction after averaging
logical states 0 and 1.  All 11 cohorts and all five methods are reported;
the forest plot separately shows the two gate-relevant differences.}
\label{tab:pnnl-cohort-results}
\small
\ifrunresultbundle
  \resizebox{\linewidth}{!}{\RunSixVerifiedPNNLCohortResultsTable}
\else
  \RunSixPending{PNNL_COHORT_RESULTS_TABLE}
\fi
\end{table}

\begin{table}[ht]
\centering
\caption{Aggregate \PNNL{} retention conditions against both predesignated
comparators, with the path bootstrap, calibration-pair sensitivity bootstrap,
exact paired sign-flip \(p\)-value, and all three gate Booleans.}
\label{tab:pnnl-comparison}
\small
\ifrunresultbundle
  \resizebox{\linewidth}{!}{\RunSixVerifiedPNNLComparisonTable}
\else
  \RunSixPending{PNNL_AGGREGATE_RETENTION_TABLE}
\fi
\end{table}

\begin{table}[ht]
\centering
\caption{\PNNL{} cohort-control map.  The anonymous C01--C11 plot labels
map to the frozen path/snapshot cohorts here; QASM equality is exposed, and
the control class is derived per cohort from the hash-bound Pittsburgh
manifest.  These constructed snapshot comparisons are not natural drift.}
\label{tab:pnnl-controls}
\small
\ifrunresultbundle
  \resizebox{\linewidth}{!}{\RunSixVerifiedPNNLCohortControlTable}
\else
  \RunSixPending{PNNL_COHORT_CONTROL_TABLE}
\fi
\end{table}

\begin{figure}[ht]
\centering
\ifrunresultbundle
  \RunSixVerifiedPNNLFigure
\else
  \fbox{\parbox[c][32mm][c]{0.90\linewidth}{\centering
  \RunSixPending{PNNL_DELAY_FOREST_FIGURE}}}
\fi
\caption{\PNNL{}/IBM constructed snapshot-boundary results.}
\label{fig:pnnl}
\end{figure}

\subsection{Resource ledger}

\begin{table}[ht]
\centering
\caption{Hash-bound resource ledger.  Timing and maximum RSS are
process-wide measurements for the joint pipeline and are not attributed to
individual methods.  All three Google joint replays, both PNNL wall-time
scopes, and PNNL bootstrap/randomization surrogate work and
eigendecompositions are reported; surrogate work is not new hardware
exposure.}
\label{tab:resources}
\small
\ifrunresultbundle
  \resizebox{\linewidth}{!}{\RunSixVerifiedResourceLedgerTable}
\else
  \RunSixPending{JOINT_PIPELINE_RESOURCE_LEDGER}
\fi
\end{table}

\subsection{Locked decision}

\begin{table}[ht]
\centering
\caption{Recomputed atomic predicates and conjunctive locked decision.}
\label{tab:gate}
\small
\ifrunresultbundle
  \resizebox{\linewidth}{!}{\RunSixVerifiedGateTable}
\else
  \RunSixPending{RECOMPUTED_GATE_DECISION_TABLE}
\fi
\end{table}

The result-contingent sentence emitted from that recomputation is:
\RunSixClaimSentence

\section{Discussion}
\label{sec:discussion}

\subsection{What a satisfied conjunctive gate would and would not mean}

If and only if both gates in Eq.~(14) pass, the supported statement is an
empirical fact: both fixed conditions were satisfied for these datasets,
implementations, endpoints, and budgets.  Because the tests compare specific
implementations at fixed endpoints, passing does not establish superiority
over same-feature or same-parity logistic or threshold classes, a correct
likelihood or oracle rule, or a general algorithmic advantage.  It also does
not establish universal sample efficiency, scaling advantage, or quantum
acceleration.  If either gate fails, the benchmark remains informative by
identifying the simpler comparator that matched or exceeded the composite
without changing the success criterion.

\subsection{Symmetry and physics language}

Sector projections are classical representation bookkeeping on observed
features.  Group-symmetric quantum hypothesis testing
~\cite{hiai2009symmetry} and symmetry-resolved entanglement
~\cite{goldstein2018symmetry} are related mathematical contexts, not Run 6
outputs.  The primary data are a one-dimensional repetition-code syndrome
archive; they do not constitute a Wilson-loop measurement, a new
topological-order theorem, or evidence about string theory or holography.

\subsection{Limitations}

The primary hardware evidence contains one author-identified event cluster.
Its onset is approximate, and archived shot order is not a calibrated
wall-clock cadence.  The paired reference reservoir is historical rather
than physically simultaneous, so exact exchangeability is supplied only by
the artificial randomization audit.  \PNNL{} comparisons are constructed
snapshot boundaries and some confound circuit with hardware changes.  The
same 300-coordinate bank is broad but finite; no monitor can detect a change
absent from its accessible syndrome sigma-field.  The consumer validator was
repaired only after detector values had been produced and numeric diagnostics
had been exposed.  Although the repair was source/schema-derived, left the
detector bundle unchanged, and preceded outcome and held-\PNNL{} access, this
precludes describing the complete execution as detector-blind or pristine
end-to-end preregistration.  Finally, the full spectral
branch costs repeated eigendecompositions, so any computational benefit must
be measured rather than inferred from parameter count.

\subsection{Practical next experiment}

The most consequential follow-up is an online surface-code experiment with
prospectively timestamped calibration interventions and synchronized
decoder likelihoods.  A resource-matched protocol should compare \SPace{}
with covariance/CUSUM, detector likelihood, adaptive decoder/noise trackers,
and predeclared gauge-invariant diagnostics at identical shots, qubit-cycle
exposure, and intervention cost.  Classical-shadow e-detection becomes a
valid comparator only if its randomized acquisition settings are collected
as part of that separate experiment.

\section{Conclusion}
\label{sec:conclusion}

This work specifies an auditable route from past-only maximum-difference
features to sequential monitoring of fixed QEC syndrome records.  Its exact
content is deliberately narrow: bounded predictable factors under a
declared conditional design, fixed mixtures at the complete-shot
experimental unit, and standard lifetime/ARL guarantees used without
interchange.  The result-contingent conclusion is:
\RunSixClaimSentence\ The
benchmark is designed so that a strong rate, covariance, likelihood, or
same-feature comparator remains a scientific result rather than a baseline
to be removed after inspection.

\section*{Data and Code Availability}

The Google and \PNNL{} data are public deposits
~\cite{google2022data,stein2026dataset}.  The accompanying repository stores
the initial preregistration, original freeze ratification, separately
ratified post-detector repair, detector/outcome separation, truthful access
record, environment lock, source hashes, and machine-readable resource
ledgers.  Final manuscript values and figures are accepted only through the
dual-provenance, hash-bound publication bundle.

\section*{Acknowledgments}

The author thanks the creators of the public QEC datasets and the developers
of the statistical and QEC tools on which this reproducible benchmark
depends.  This working paper makes no claim of endorsement by those groups.

\appendix

\section{Frozen Hyperparameters}
\label{app:hyperparameters}

\begin{table}[ht]
\centering
\caption{Selected locked parameters.}
\small
\begin{tabularx}{\linewidth}{@{}lX@{}}
\toprule
Object & Frozen value\\
\midrule
M3 learning rates & \(0.001,0.01,0.1\); L2 \(10^{-4}\), no intercept\\
M4 half-lives & \(4,16,64,256\) role updates\\
M4 supports & \(k=1,4,16,64\)\\
M5 half-lives & \(4,16,64\) role updates\\
M5 effects & stride 8; \(\tau_{\rm eig}=10^{-10}\); deterministic anchored
rank one and full positive eigenspace\\
Bet fractions & \(0.1,0.3,0.6,0.9\)\\
Proper-prior threshold & \(\alpha=0.01\), crossing at \(100\)\\
Empirical cycle budget & at most 2 notifications in 255,000 opportunities\\
Risk budgets & top \(2,20,200\) of 20,000 monitor shots\\
Google randomization & 256 predeclared complete-shot swap streams\\
Google threshold uncertainty & 2,000 shot-block draws, block length 128\\
PNNL stream threshold & 4,096 shot-block draws, block length 32\\
PNNL paired-swap audit & 256 seeds; 22 path--state episodes; five methods\\
PNNL uncertainty & 10,000 path draws; 10,000 cluster sensitivity draws;
all \(2^{11}\) paired signs\\
\bottomrule
\end{tabularx}
\end{table}

\section{Proof Detail for the Proper-Prior Process}
\label{app:prior}

Let \(\rho_u=1/H\) for \(u=1,\ldots,H\) and
\(\rho_{>u}=(H-u)/H\).  From Eq.~(8) and conditional e-validity,
\[
\E_0[A_{u,r,c}\mid\cF_{u-1}]
\le A_{u-1,r,c}+\rho_u.
\]
After fixed mixing,
\[
\E_0[E_u\mid\cF_{u-1}]
\le \rho_{>u}
+\frac1{51}\sum_{r,c}\pi_c(A_{u-1,r,c}+\rho_u)
=E_{u-1}.
\]
The initial tail mass gives \(E_0=1\); Ville's inequality yields
Eq.~(12).  Replacing the proper start prior by an unnormalized SR clock
changes the guarantee and cannot retain the same lifetime interpretation.

\section{Result-Insertion Checklist}
\label{app:checklist}

Before a submission build, all of the following must hold:

\begin{enumerate}
  \item every keyed pending marker is absent from the compiled PDF;
  \item the hash-bound publication bundle is independently regenerated from
  all eight declared evidence inputs and matches every included byte;
  \item the immutable detector binds the original pre-access ratification,
  while every downstream arm and the final decision bind the separately
  ratified repair;
  \item the repair access record discloses diagnostic exposure and rejects
  detector-blind or pristine end-to-end preregistration language;
  \item the abstract calls the Google onset approximate and \PNNL{} a
  constructed snapshot boundary;
  \item exact e-validity is confined to the randomized complete-pair design;
  \item all 51 roles remain one formal time step;
  \item e-process lifetime control and SR ARL are not conflated;
  \item paired \DFR{} and same-feature logistic both appear in the gate;
  \item M0C, M1, and M2 are reported regardless of outcome;
  \item the complete threshold/frontier and locked uncertainty tables are
  included, including the descriptive 2,000-replicate threshold bootstrap;
  \item PNNL pre-false-alarm, miss, macro-delay, and both comparator
  retention conditions are included without selecting paths, together with
  the complete paired-swap implementation-audit summary;
  \item the recorded chain shows zero completed randomization replicates and
  no outcome or held-\PNNL{} access before repair, with no detector rerun;
  \item uncertainty is clustered by shot/event or path as declared;
  \item a tied or failed gate produces the locked negative conclusion; and
  \item no quantum speedup, universal sample-efficiency, topological,
  string-theory, or holographic claim appears.
\end{enumerate}

\ifrunrevtex
  \bibliographystyle{apsrev4-2}
\else
  \bibliographystyle{unsrtnat}
\fi
\bibliography{references}

\end{document}
